%% 第31回編集委員会（物理第5巻）：組 原W6
%% p5-11.tex  原-11　質量はエネルギー：相対論と原子核（★3〜5・3段）
%% 中身は計画書 tools/phys5-plan.md の B「原-11」。公式の持ち主は A（A3〔中性子〕・A4・A22〔導く〕・A30〜A33・A36・A38・A39・A66〜A68）。
%%   ★は計画書 B の表（N-1）の 3〜5 にした（series_phys.py も 3〜5 に直してある：2aeb619）。
%% 数値はすべて Python で計算した：checks/物理第5巻/原-11/calc_p5_11.py（出力 calc_p5_11_out.txt）。
%%   原子の質量は同じフォルダの nist_masses_20261111.json（NIST の表を 11月11日に保存したもの）から読む。
%% 前の巻・数学の本との約束（計画書 A-6・C-6-2）
%%   ・Q1（p1-01:806 速度の合成は相対論で）→ ⑦・例題7（0.6c＋0.6c＝0.882c。力-1 の 0.5c＋0.5c＝0.8c も）
%%   ・Q2（p1-02:909 運動量が mv でなくなる）→ ⑧。Q3（p1-03:706 衝突は相対論の運動量で）→ ⑧のコンプトン散乱
%%   ・Q20（p2-16:335 太陽は核融合、ITER 50 MW→500 MW）→ はてな？・⑪・例題4
%%   ・Q31（p3-08:602 光のドップラー効果の正しい式）→ ⑨・例題11（0.905、音の式 0.909、近似 0.900）
%%   ・Q35（p3-13:1086 マイケルソン・モーリー 1887年）→ ①
%%   ・Q37（p3-14:347・p1-01:812 GPS の時計 約38 μs/日）→ ⑦・例題6（−7.2・＋45.7・＋38.5 μs/日、高さ 20200 km は波-14 と同じ）
%%   ・Q40（p4-01:707・p4-02:879 1 kV の電子は光の約6 %）→ ⑧・例題8（0.062、100 kV で 0.55）
%%   ・Q42（p4-01:1291 陽子どうし 1 fm で 230 N）→ ⑩・例題12（8.99×10^9 を使う）
%%   ・Q52（p4-06:847 加速器の相対論）→ ⑧・工学ではこう使う。Q57（p4-14:521・1016 電場と磁場の入れかわり、1905年の論文）→ ①⑫
%%   ・A22（コンプトンの式：読むのは原-2 ⑦、導くのはこの章 ⑧）→ ⑧・例題10
%%   ・数学 2-8（√）・4-4（三平方：光の時計）・5-17（近似 √(1−x)）→ ⑦⑧の思い出そう
%% 記号（この章）：光速 c、速さ v、β＝v/c（この章だけ。β線は「β線」とことばで書く）、ローレンツ因子 γ（γ線は「γ線」とことばで）、
%%   質量 m、静止エネルギー mc^2、運動量 p、エネルギー E、運動エネルギー K、原子番号 Z・中性子数 N・質量数 A（この章では A は質量数だけ。
%%   放射能は記号にしない）、結合エネルギー B（この章では磁場は出さない）、質量欠損 Δm、原子質量単位 u（立体）、Q 値 Q、
%%   原子核の半径 R・r_0、振動数 f（N-5）、波長 λ（この章の λ は光の波長だけ。崩壊定数は原-12 から）。
%% 出典（11月11日に確かめた。Wikipedia は使っていない）
%%   ・定数：CODATA 2022（NIST allascii.txt。計画書 D4）。本文の丸め c＝3.00×10^8、1 u＝931.5 MeV、m_e c^2＝0.511 MeV。
%%   ・原子の質量：NIST "Atomic Weights and Isotopic Compositions"（AME2020 にもとづく）
%%       https://physics.nist.gov/cgi-bin/Compositions/stand_alone.pl?ele=He&ascii=ascii2&isotype=all（ele を元素ごとに変えて読んだ）。
%%   ・太陽の明るさ 3.828×10^26 W、太陽の GM（名目値）：IAU 2015 Resolution B3（arXiv:1510.07674。計画書 D46）。
%%   ・ミューオンの平均寿命 2.1969803 μs：MuLan, arXiv:1211.0960 の要旨（本文は 2.2 μs に丸めた）。
%%   ・GPS の軌道半径 26560 km（高さ 20200 km）：米国の GPS の仕様書 IS-GPS-200N（https://www.gps.gov/sites/default/files/2025-07/IS-GPS-200N.pdf）の
%%       軌道の長半径の基準値 A_REF＝26,559,710 m と、gps.gov の "Space Segment"（https://www.gps.gov/space-segment）の
%%       "an altitude of approximately 20,200 km"（出典の組が 11月11日に curl で読んだ。はじめは ESA Navipedia の抜粋）。波-14 の 20200 km と同じ。
%%   ・水の生成エンタルピー −285.83 kJ/mol：NIST Chemistry WebBook（webbook.nist.gov、ID=C7732185）。
%%   ・ノーベル賞（nobelprize.org で読んだ）：1921 アインシュタイン、1935 チャドウィック（「1932年に中性の粒子であることを示した」）、
%%       1949 湯川（「1935年に論文」）、1950 パウエル（中間子）、1951 コッククロフトとウォルトン（「1932年にリチウムとホウ素を陽子で変換」）、
%%       2002 小柴昌俊、2015 梶田隆章（計画書 D33・D36）。
%%   ・マイケルソン・モーリー 1887年：波-13（p3-13:1086）と同じ。1905年の論文の題：電-14（p4-14:521）と同じ。
%%   ・ニュートリノの質量の上限 0.45 eV（KATRIN）：arXiv:2406.13516（計画書 D44）。
%%   ・太陽の中心の温度 約1500万 K：文部科学省の資料（mext.go.jp/stw の「太陽のことば」PDF）と日本物理学会誌 71巻4号の記事の2つで照らし合わせた（検索の抜粋）。
%%   ・π中間子（π±）の質量 139.57039±0.00018 MeV：PDG 2025 の Listings（https://pdg.lbl.gov/2025/listings/rpp2025-list-pi-plus-minus.pdf を
%%       curl で読んだ。出典の組、11月11日）。本文は「約140 MeV」（139.57 MeV）。
%% 〔要確認〕（確かめられなかった値）
%%   ・半経験的質量公式の係数 15.8・18.3・0.714・23.2・12.0 MeV：教科書ごとに少しずつちがう係数の一例（計画書と同じ）。本文に「一例」と書いた。
%%   ・原子核の半径の係数 r_0＝1.2 fm：教科書の目安の値（1.2〜1.25 fm）。本文に「約」と書いた。
%%   ・地球の GM＝3.986×10^14 m^3/s^2（計画書と同じ）、石炭の発熱量 3×10^7 J/kg（はてな？の見積もりの仮の値。本文に「とします」）。
%% 絵の部品の名前：txl…（原-11 → l）
\makeatletter
% 陽子（＋の印）と中性子（印なし）  \txlp{(x,y)}{半径}  \txln{(x,y)}{半径}
\newcommand\txlp[2]{\filldraw[fill=bkRed!35, draw=bkRed!70!black, line width=0.5pt] #1 circle (#2);
  \draw[bkRed!70!black, line width=0.6pt] #1 ++(-0.45*#2,0) -- ++(0.9*#2,0) #1 ++(0,-0.45*#2) -- ++(0,0.9*#2);}
\newcommand\txln[2]{\filldraw[fill=bkGray!25, draw=bkGray!80!black, line width=0.5pt] #1 circle (#2);}
% 鏡（光の時計）  \txlmirror{(x,y)}{はば}
\newcommand\txlmirror[2]{\fill[bkBlue!25] ($#1+(-#2/2,-0.6)$) rectangle ($#1+(#2/2,0.6)$);
  \draw[bkInk, line width=0.7pt] ($#1+(-#2/2,-0.6)$) rectangle ($#1+(#2/2,0.6)$);}
\tikzset{txl lab/.style={font=\sffamily\footnotesize, text=bkInk}, txl arr/.style={line width=1pt, -{Stealth[length=2.4mm]}},
  txl dim/.style={bkGray, line width=0.5pt, {Stealth[length=1.8mm]}-{Stealth[length=1.8mm]}},
  txl light/.style={bkOrange!85!black, line width=1.1pt, -{Stealth[length=2.2mm]}}}
\makeatother

\phchapter{11}{質量はエネルギー：\\相対論と原子核}{太陽は、何をエネルギーのもとにして光っている？}{\bklevelrange{3}{5}}{質量はエネルギー：相対論と原子核}
\bkmichishirube{2-8,4-4,5-17,3-12,力-1,力-3,波-8,波-13,電-1,電-14,原-1,原-2,原-6}{原-12,原-13,原-15,原-16}

\begin{goal}
\begin{itemize}[label={\color{bkTheme}●}, leftmargin=1.6\zw, itemsep=0pt, topsep=1pt]
\item \bklevel{3}\enspace 原子番号・質量数・同位体ということばで、原子核のなかみを説明できる。
\item \bklevel{3}\enspace 原子の質量の表から、結合エネルギーと核反応の Q 値を計算できる。
\item \bklevel{4}\enspace 光の時計から時間のおくれを導き、GPS の時計のずれを計算できる。
\item \bklevel{4}\enspace $E^2=(pc)^2+(mc^2)^2$ から、$E=mc^2$ とコンプトンの式を導ける。
\item \bklevel{5}\enspace 半経験的質量公式で、結合エネルギーの曲線の形を説明できる。
\end{itemize}
\end{goal}

\phkoushiki{%
\phk{原子核は陽子と中性子、同位体}{中学理科}{ph:原-11:sec3}
\phk{原子番号・質量数 $A=Z+N$、u}{高校物理}{ph:原-11:nuc}
\phk{質量とエネルギー $E=mc^2$}{高校物理}{ph:原-11:emc2}
\phk{結合エネルギー $\Delta E=\Delta m\,c^2$}{高校物理}{ph:原-11:sec4}
\phk{核反応式と、核反応で出るエネルギー}{高校物理}{ph:原-11:sec6}
\phk{核力は、短い距離で強く引き合う}{高校物理}{ph:原-11:sec10}
\phk{軽い原子核が合わさる核融合}{高校物理}{ph:原-11:sec11}
\phk{素粒子（クォークとレプトン）}{高校物理}{ph:原-11:soryushi}
}

\begin{hatena}
\noindent\begin{minipage}[c]{0.56\linewidth}\raggedright
太陽は、地球から1億5千万 km もはなれているのに、地球の大気の外側で、太陽に向けた 1 m$^2$ の面に毎秒 \qty{1361}{J} もの光を届けています。太陽が四方八方へ出す光を全部合わせると、1秒に $\qty{3.83e26}{J}$ です。

もし太陽が、石炭のように燃えて光っているとしたら、どうなるでしょう。太陽の質量は約 $\qty{2.0e30}{kg}$ です。石炭 1 kg が燃えて出す熱を $\qty{3e7}{J}$ とすると、太陽全部が燃えても、約5000年で燃えつきてしまいます。でも、太陽は約46億年も光り続けています。太陽は、何をエネルギーのもとにしているのでしょう。
\end{minipage}\hfill
\begin{minipage}[c]{0.41\linewidth}\centering
\begin{tikzpicture}[x=0.85mm, y=0.85mm]
  \fill[bkYellow!35] (24,24) circle (21);
  \fill[bkOrange!35] (24,24) circle (12);
  \fill[bkRed!30] (24,24) circle (5);
  \draw[bkOrange!80!black, line width=0.7pt] (24,24) circle (21);
  \foreach \a in {0,30,...,330} {\draw[bkOrange!80!black, line width=0.8pt] (24,24) ++(\a:23) -- ++(\a:4);}
  \node[txl lab, align=center] at (24,24) {中心\\？};
  \node[txl lab, anchor=north, align=center] at (24,-5) {光り続けて約46億年};
\end{tikzpicture}
\end{minipage}
\end{hatena}

\noindent 第4部「分かれる」では、原子の中心にある\textbf{原子核}の物理を学びます。原子核は、原子の大きさの数万分の1しかない小さな粒ですが、原子の質量のほとんどを持ち、化学反応の何百万倍ものエネルギーをしまっています。そのエネルギーを数える道具が、アインシュタインの\textbf{相対性理論}と、そこから出てくる式 $E=mc^2$ です。この章では、相対性理論の入口と、原子核の質量とエネルギーの関係を学び、はてな？に答えます。続く章で、こわれる原子核（原-12）、放射線と物質（原-13・14）、原子炉（原-15）へ進みます。

\begin{column}{安全：第4部を読む前に}
第4部（原-11〜16）では、放射線や、放射線を出す物質（放射性物質）が出てきます。この本の「やってみよう」では、放射線を出す物や装置は使いません。
\begin{enumerate}[label={\bkcircled{\arabic*}}, leftmargin=2.4\zw, itemsep=0pt, topsep=2pt]
\item 放射線を出すかもしれない物（ある種の鉱物、古い夜光塗料の時計の文字盤など）を、自分で集めたり、こわしたり、体に近づけたりしない。
\item 放射線は、学校・自治体・科学館の測定器で、先生や係の人といっしょにはかる。自治体が公開している空間線量率のデータを見るのもよい（原-14）。霧箱などは、科学館などで係の人がいる所で見る。
\end{enumerate}
放射線の影響は「あるか、ないか」ではなく、受けた\textbf{量}で決まります（原-14）。
\end{column}

\begin{yattemiyou}[原子の質量を足し算してみよう]
\textbf{用意するもの}：電卓、表\ref{ph:原-11:tab-yattemiyou}\par\vspace{2pt}
\noindent\begin{minipage}[c]{0.56\linewidth}
\begin{enumerate}[label={\bkcircled[bkGreen]{\arabic*}}, leftmargin=2.4\zw, itemsep=1pt, topsep=0pt]
\item ヘリウム原子は、陽子2個・中性子2個・電子2個でできています。水素原子（陽子1個と電子1個）2個と、中性子2個の質量を足しましょう。
\item その和を、ヘリウム原子の質量と比べましょう。どちらが重いですか。
\item 炭素12（陽子・中性子6個ずつ）、酸素16（8個ずつ）でも、足して比べましょう。
\end{enumerate}
\end{minipage}\hfill
\begin{minipage}[c]{0.41\linewidth}\centering
\begin{bktable}{colspec={Q[l,m,wd=22mm] Q[r,m,wd=24mm]}, rowsep=0.3pt, colsep=3pt, row{1}={bg=bkTheme!10}, cells={font=\small}}
もの & 質量〔u〕 \\
水素原子 ${}^1$H & 1.007825 \\
中性子 & 1.008665 \\
${}^4$He 原子 & 4.002603 \\
${}^{12}$C 原子 & 12（ちょうど） \\
${}^{16}$O 原子 & 15.994915 \\
\end{bktable}
\bktabcaption{NIST の表の質量（u は③）}\label{ph:原-11:tab-yattemiyou}
\end{minipage}
\end{yattemiyou}

\noindent やってみようの答え：ヘリウムの場合、$2\times1.007825+2\times1.008665=4.032980$ u です。ところが、本物のヘリウム原子は 4.002603 u で、\textbf{0.030377 u だけ軽い}のです。炭素12では、和が 12.098940 u なのに本物は 12 u（0.098940 u 軽い）、酸素16では、和が 16.131920 u なのに本物は 15.994915 u（0.137005 u 軽い）でした。陽子と中性子が集まって原子核になると、ばらばらのときより、いつも少し軽くなります。消えた質量はどこへ行ったのでしょう。この章の答えは「\textbf{エネルギーになって出ていった}」です（④）。

\begin{omoidasou}
\bkrecall{原-1}{原子は、中心の小さな原子核（正の電気を持ち、質量のほとんどが集まる）と、周りの電子でできている。}
\bkrecall{電-1}{距離 $r$ はなれた電荷 $q_1$・$q_2$ の間のクーロン力は $F=kq_1q_2/r^2$（$k=\qty{8.99e9}{N.m^2/C^2}$）。}
\bkrecall{電-2}{電子を 1 V で加速したときのエネルギーが $\qty{1}{eV}=\qty{1.60e-19}{J}$（$\qty{1}{MeV}=10^6$ eV）。}
\bkrecall{原-2}{光子のエネルギーは $E=hf$、運動量は $p=h/\lambda=E/c$（$f$ は振動数。高校の教科書の $h\nu$ と同じ）。}
\bkrecall{4-4}{直角三角形の辺 $a$・$b$ と斜辺 $c$ には $a^2+b^2=c^2$（三平方の定理）。}
\bkrecall{5-17}{$x$ が小さいとき $1/\sqrt{1-x}\approx1+x/2$（1次近似）。}
\end{omoidasou}

\noindent この章の記号：$c$ は光の速さ、$v$ は物の速さです。光の速さと比べた速さ $\beta=v/c$ と、ローレンツ因子 $\gamma$（⑦）は、この章だけの記号です。放射線のβ線・γ線とは別のものなので、放射線のほうは「β線」「γ線」とことばで書きます。$A$ は質量数、$B$ は結合エネルギー、$Q$ は核反応で出るエネルギーです。

%% ======================== 物理の道具 ========================
\Needspace{10\baselineskip}\bksection{doguu}{物理の道具}

\phstage{me}{光の速さと、質量のエネルギー}{3}

\bksubsection{1}{光の速さは、だれから見ても同じ}
\label{ph:原-11:sec1}
時速 \qty{100}{km} で走る電車の中で、前へ時速 \qty{30}{km} でボールを投げると、地面から見たボールの速さは時速 \qty{130}{km} です（力-1）。では、光の速さの半分で飛ぶロケットから、前へ懐中電灯の光を出したら、地面から見た光の速さは、光の速さの1.5倍になるでしょうか。

答えは「いいえ」です。地面の人がはかっても、ロケットの人がはかっても、光の速さは同じ $c$ です（図\ref{ph:原-11:fig-rocket}）。
\begin{bkfigure}
\begin{tikzpicture}[x=0.85mm, y=0.85mm]
  \draw[bkGray, line width=0.6pt] (0,2) -- (52,2);
  \fill[bkBlue!25] (4,8) -- (20,8) -- (24,11) -- (20,14) -- (4,14) -- cycle;
  \draw[bkInk, line width=0.6pt] (4,8) -- (20,8) -- (24,11) -- (20,14) -- (4,14) -- cycle;
  \draw[bkInk, txl arr] (8,18) -- (18,18);
  \node[txl lab, anchor=south] at (13,19) {$0.5c$};
  \draw[txl light] (26,11) -- (50,11);
  \node[txl lab, anchor=south] at (38,12) {光};
  \node[txl lab, anchor=north, align=center] at (26,0) {地面から見ても $c$\\（$1.5c$ ではない）};
\end{tikzpicture}
\bkfigcaption{光の速さの半分で飛ぶロケットから出した光も、地面から見た速さは $c$。}\label{ph:原-11:fig-rocket}
\end{bkfigure}
このことは、1887年のマイケルソンとモーリーの実験で確かめられました（物理シリーズの『波動』（波-13）の干渉計）。地球は太陽の周りを秒速約30 km で回っているので、光の速さが「地球の動きと同じ向き」と「逆向き」で違えば、干渉じまが動くはずでした。でも、しまは動きませんでした。

1905年、アインシュタインは、次の2つを出発点にして、時間と空間の考え方を作り直しました。これが\textbf{特殊相対性理論}です。\bkidx{とくしゅそうたいせいりろん}{特殊相対性理論}
\begin{enumerate}[label={\bkcircled{\arabic*}}, leftmargin=2.4\zw, itemsep=1pt, topsep=2pt]
\item 一定の速さでまっすぐ動く人から見ても、止まっている人から見ても、物理の法則は同じ形をしている（相対性原理）。
\item 光の速さ $c$ は、光を出す物の速さにも、見る人の速さにもよらず、同じ（光速の不変）。
\end{enumerate}
この論文の題は「動いている物体の電気力学について」でした。電-14 で見たように、止まった人には電場だけに見えるものが、動く人には磁場もふくんで見えます。電場と磁場の入れかわりを調べることが、相対性理論の入口だったのです。

いまでは、光の速さは \qty{299792458}{m/s} と決められていて、1 m は「光が $1/\num{299792458}$ 秒に進む距離」です。この本では、ふつう $c=\qty{3.00e8}{m/s}$ を使います。

\begin{phmijika}[相対性理論は、身近な道具の中にある]
カーナビやスマートフォンの位置を決める GPS の衛星の時計は、相対性理論で計算したぶんだけ補正されています（⑦・例題6）。病院の放射線治療の加速器や、PET 検査（⑪）も、相対性理論なしには設計できません。光の速さに近い物は、わたしたちの身の回りにも、たくさんあるのです。
\end{phmijika}

\Needspace{8\baselineskip}\bksubsection{2}{$E=mc^2$：質量はエネルギーのかたまり}
\label{ph:原-11:sec2}
相対性理論から出てくる、いちばん有名な式が
\pheq{E=mc^2}{J ＝ kg × (m/s)$^2$}
です（⑧で導きます）。止まっている物も、質量 $m$ に応じて $mc^2$ のエネルギーを持っています。これを\textbf{静止エネルギー}といいます。\bkidx{せいしえねるぎー}{静止エネルギー}逆に、物がエネルギーを出すと、そのぶんだけ質量が減ります。

\begin{phyomi}[$E=mc^2$]
「質量 $m$〔kg〕の物は、$m$ に光の速さの2乗をかけただけのエネルギー〔J〕を持つ」と読みます。$c^2=(3.00\times10^8)^2=9.0\times10^{16}$ はとても大きな数なので、ほんの少しの質量が、とても大きなエネルギーになります。たとえば \qty{1}{g} の質量は
\[ E=1\times10^{-3}\times9.0\times10^{16}=\qty{9.0e13}{J} \]
で、\qty{100}{W} の電球を約3万年つけ続けられるエネルギーです。
\end{phyomi}

\begin{phbaisuru}{質量 $m$}
\noindent\begin{minipage}[c]{0.5\linewidth}\raggedright
$E$ は $m$ に比例するので、質量が2倍なら静止エネルギーも2倍です。グラフは原点を通る直線で、傾きが $c^2$ です。傾きがとても大きいので、横軸を\textbf{マイクログラム}（\si{\micro g}）にしても、縦軸は何千万 J にもなります。
\end{minipage}\hfill
\begin{minipage}[c]{0.46\linewidth}\centering
\begin{tikzpicture}
\begin{axis}[bkplot, width=44mm, height=34mm, xmin=0, xmax=2.4, ymin=0, ymax=2.4e8, xtick={1,2}, ytick={0.9e8,1.8e8}, yticklabels={$9\times10^7$,$1.8\times10^8$}, xlabel={$m$（\si{\micro g}）}, ylabel={$E$（J）}]
  \addplot[bkcurve, domain=0:2.2] {9e7*x};
  \draw[bkGray, dashed, line width=0.5pt] (axis cs:1,0) -- (axis cs:1,9e7) -- (axis cs:0,9e7);
  \draw[bkGray, dashed, line width=0.5pt] (axis cs:2,0) -- (axis cs:2,1.8e8) -- (axis cs:0,1.8e8);
\end{axis}
\end{tikzpicture}
\end{minipage}
\end{phbaisuru}

\phboxhead{化学の反応と、原子核の反応}
では、なぜ、ふだんは質量が減ることに気づかないのでしょう。水素が燃えて水ができるとき、水の分子1個あたりに出る熱は約 \qty{3.0}{eV} です（水 1 mol あたり \qty{286}{kJ}）。水の分子1個の静止エネルギーは約 $\qty{1.7e4}{MeV}$ なので、減る質量は全体の約 $1.8\times10^{-10}$、つまり100億分の2ほどで、どんな天びんでもはかれません。お湯 1 kg を \qty{80}{K} 温めると、お湯の質量は $\qty{3.7e-12}{kg}$ だけ増えますが、これもはかれません。

ところが、原子核の反応では、1回で何 MeV（何百万 eV）ものエネルギーが出ます。ウランの原子核が2つに分かれる（核分裂）と、1回で約 \qty{200}{MeV}、化学の反応（水素が燃えるとき約 \qty{3}{eV}）の約7千万倍です（図\ref{ph:原-11:fig-chem}）。減る質量も全体の約0.1\,\% になり、原子の質量の表から、はっきり読みとれます。

\begin{bkfigure}
\begin{tikzpicture}
\begin{axis}[bkplot, width=118mm, height=40mm, xmode=log, xmin=1, xmax=1e9, ymin=0.4, ymax=3.6, ytick={1,2,3}, yticklabels={水素が燃える,D と T の核融合,ウランの核分裂}, xtick={1,1e3,1e6,1e9}, xticklabels={1 eV,1 keV,1 MeV,1 GeV}, axis lines=left, grid=none]
  \foreach \v/\y in {2.96/1,17.59e6/2,200e6/3} {\edef\tmp{\noexpand\filldraw[fill=bkBlue!45, draw=bkInk, line width=0.5pt] (axis cs:1,\y-0.22) rectangle (axis cs:\v,\y+0.22);}\tmp}
\end{axis}
\end{tikzpicture}
\bkfigcaption{1回の反応で出るエネルギー（横の目もりは対数で、1目もりで1000倍）。水素が燃えるのは約 3 eV、重水素（D）と三重水素（T）の核融合は 17.6 MeV（例題3）、ウランの核分裂は約 200 MeV（原-15）。}\label{ph:原-11:fig-chem}
\end{bkfigure}

\Needspace{12\baselineskip}\phstage{te}{質量から、エネルギーを数える}{3}

\bksubsection{3}{原子核のなかみと同位体}
\label{ph:原-11:sec3}
原子核は、正の電気を持つ\textbf{陽子}と、電気を持たない\textbf{中性子}でできています。陽子と中性子をまとめて\textbf{核子}といいます。\bkidx{かくし}{核子}陽子の数を\textbf{原子番号} $Z$、中性子の数を $N$、核子の数を\textbf{質量数} $A$ といい、
\pheq{A=Z+N}{個 ＝ 個 ＋ 個}
\label{ph:原-11:nuc}です。\bkidx{しつりょうすう}{質量数}原子番号は元素の種類（化学の性質）を決めます。原子核は、元素の記号の左上に質量数、左下に原子番号を書いて ${}^{4}_{2}\mathrm{He}$ のように表します。原子番号は元素の記号で決まるので、左下は書かずに ${}^{4}\mathrm{He}$ や「ヘリウム4」と書くことも多いです。

同じ元素（同じ $Z$）でも、中性子の数が違う原子核があります。これを\textbf{同位体}といいます。\bkidx{どういたい}{同位体}水素には、ふつうの水素 ${}^1$H のほかに、中性子を1個持つ重水素 ${}^2$H（D とも書く）と、2個持つ三重水素 ${}^3$H（トリチウム、T）があります。炭素にも ${}^{12}$C・${}^{13}$C・${}^{14}$C などがあります。同位体は化学の性質がほとんど同じで、質量だけが違います（図\ref{ph:原-11:fig-iso}）。
\begin{bkfigure}
\begin{tikzpicture}[x=1mm, y=1mm]
  \txlp{(6,10)}{2.4}
  \node[txl lab, anchor=north] at (6,4) {${}^1$H};
  \txln{(26.6,10)}{2.4} \txlp{(22,10)}{2.4}
  \node[txl lab, anchor=north] at (24,4) {${}^2$H（D）};
  \txln{(46.4,11.8)}{2.4} \txln{(44.2,8)}{2.4} \txlp{(42,11.8)}{2.4}
  \node[txl lab, anchor=north] at (44,4) {${}^3$H（T）};
\end{tikzpicture}
\bkfigcaption{水素の3つの同位体（＋の印が陽子、印のない丸が中性子）}\label{ph:原-11:fig-iso}
\end{bkfigure}
中性子は、電気を持たないので見つけにくく、最後に見つかった核子です。1930年、ベリリウムに α粒子（ヘリウムの原子核）を当てると、物をよく通りぬける何かが出てくることがわかりました。1932年、イギリスのチャドウィックは、それが陽子とほぼ同じ質量を持つ、電気を持たない粒であることを示しました。今の書き方では
\[ {}^{9}_{4}\mathrm{Be}+{}^{4}_{2}\mathrm{He}\longrightarrow{}^{12}_{6}\mathrm{C}+{}^{1}_{0}\mathrm{n} \]
という反応です（n は中性子。⑥で読み方を学びます）。チャドウィックは、この発見で1935年のノーベル物理学賞を受けました。

\phboxhead{質量の単位 u}
原子や原子核の質量は、kg で書くと小さすぎるので、\textbf{原子質量単位} u を使います。\bkidx{げんししつりょうたんい}{原子質量単位}炭素12の原子1個の質量のちょうど12分の1を 1 u と決めます。
\[ \qty{1}{u}=\qty{1.66054e-27}{kg} \]
陽子は 1.007276 u、中性子は 1.008665 u、電子は 0.000549 u です。核子はどれも約 1 u なので、原子の質量は u で書くと、ほぼ質量数と同じ数になります。

\phboxhead{原子核の大きさ}
いろいろな原子核に電子や α粒子を当てて調べると、原子核の半径 $R$ は、質量数 $A$ の3乗根にほぼ比例することがわかります。
\pheq{R\approx r_0A^{1/3},\qquad r_0\approx\qty{1.2}{fm}}{m ＝ m}
fm（フェムトメートル）は $10^{-15}$ m です。半径が $A^{1/3}$ に比例するということは、体積 $\tfrac43\pi R^3$ が $A$ に比例する、つまり\textbf{核子1個が占める体積は、どの原子核でもほぼ同じ}ということです。原子核は、水のように「つめこまれたつぶの集まり」で、密度がどれもほぼ等しいのです（例題1）。原子の半径は約 0.1 nm（$10^{-10}$ m）なので、原子核は原子の1万分の1から10万分の1の大きさです。原子を直径 100 m の球にすると、原子核は直径 \qtyrange{2}{7}{mm} ほどの粒です。

\begin{phbaisuru}{質量数 $A$}
\noindent\begin{minipage}[c]{0.55\linewidth}\raggedright
質量数が2倍になると、体積は2倍、半径は $2^{1/3}\approx1.26$ 倍になります。半径はあまり大きくなりません。質量数 27 のアルミニウムは半径 \qty{3.6}{fm}、その約8.8倍の質量数 238 のウランでも、半径は約2倍の \qty{7.4}{fm} です。
\end{minipage}\hfill
\begin{minipage}[c]{0.42\linewidth}\centering
\begin{tikzpicture}[x=1mm, y=1mm]
  \filldraw[fill=bkOrange!25, draw=bkInk, line width=0.6pt] (8,10) circle (6);
  \filldraw[fill=bkOrange!25, draw=bkInk, line width=0.6pt] (34,10) circle (7.56);
  \node[txl lab] at (8,10) {$A$};
  \node[txl lab] at (34,10) {$2A$};
  \node[txl lab, anchor=north, align=center] at (8,2) {半径 1};
  \node[txl lab, anchor=north, align=center] at (34,1) {半径 1.26};
\end{tikzpicture}
\end{minipage}
\end{phbaisuru}

\Needspace{12\baselineskip}\bksubsection{4}{質量欠損と結合エネルギー}
\label{ph:原-11:sec4}
やってみようで見たとおり、陽子と中性子が集まってできた原子核は、ばらばらの核子の質量の和より軽くなります。この差を\textbf{質量欠損}といいます。\bkidx{しつりょうけっそん}{質量欠損}ヘリウム4では
\[ \Delta m=(2\times1.007825+2\times1.008665)-4.002603=\qty{0.030377}{u} \]
です。ここで、陽子の代わりに水素原子の質量、原子核の代わりにヘリウム原子の質量を使いました。どちらにも電子が2個ずつ入っているので、引き算で電子の分は消えます（原子の中の電子をつなぎとめるエネルギーは eV の大きさで、MeV の計算では無視できます）。

この質量は、どこへ行ったのでしょう。核子が集まって原子核になるとき、エネルギーが γ線などになって外へ出ていきます。出ていったエネルギーのぶんだけ、$E=mc^2$ で質量が減ったのです。\textbf{逆に、原子核をばらばらの核子に分けるには、同じだけのエネルギーを外から加えなければなりません}。このエネルギーを\textbf{結合エネルギー} $B$ といいます。\bkidx{けつごうえねるぎー}{結合エネルギー}
\pheq{B=\Delta m\,c^2=\bigl(Zm_{\mathrm H}+Nm_{\mathrm n}-M\bigr)c^2}{J ＝ kg × (m/s)$^2$}
$m_{\mathrm H}$ は水素原子、$m_{\mathrm n}$ は中性子、$M$ は調べる原子の質量です。

\phboxhead{1 u は 931.5 MeV}
1 u の質量を、静止エネルギーにすると
\[ \qty{1.66054e-27}{kg}\times(\qty{2.998e8}{m/s})^2=\qty{1.4924e-10}{J}=\qty{931.5}{MeV} \]
です。だから、u で書いた質量の差に 931.5 をかけると、MeV のエネルギーになります。ヘリウム4の結合エネルギーは
\[ B=0.030377\times931.5=\qty{28.3}{MeV} \]
で、核子1個あたり $28.3\div4=\qty{7.07}{MeV}$ です。

\begin{phflow}[$E=mc^2$ から、結合エネルギーへ]
\phflowrow{$E=mc^2$}{くっついた状態と\\ばらばらの状態の差}{$\Delta E=\Delta m\,c^2$}
\phflowrow{$\Delta E=\Delta m\,c^2$}{u と MeV で\\数える}{$B=931.5\times\Delta m$\\（$B$ は MeV、$\Delta m$ は u）}
\end{phflow}

\begin{phyomi}[$B=\Delta m\,c^2$]
「原子核の結合エネルギーは、核子がばらばらのときより軽くなった質量に $c^2$ をかけたもの」と読みます。結合エネルギーが大きいほど、原子核は\textbf{軽く}、こわれにくいのです。重くなるのではありません（まちがえやすいところ）。
\end{phyomi}

\begin{phmijika}[くっついた磁石]
2つの強い磁石がパチンとくっつくと、音が出て、ぶつかった所が少し温まります。くっつくときに、エネルギーが外へ出たのです。はがすには、出ていったぶんの仕事を、手でしなければなりません。原子核も同じで、核子がくっつくとき γ線などでエネルギーを出し、はがすにはそのぶんのエネルギーが要ります。たとえば、中性子が水素の原子核（陽子）につかまって重水素になるとき、\qty{2.22}{MeV} の γ線が出ます。これがちょうど重水素の結合エネルギーです。
\end{phmijika}

\Needspace{14\baselineskip}\bksubsection{5}{結合エネルギーの曲線}
\label{ph:原-11:sec5}
NIST の原子の質量の表から、いろいろな原子核の「核子1個あたりの結合エネルギー」$B/A$ を計算して、質量数 $A$ を横にとって並べると、図\ref{ph:原-11:fig-curve}のようになります（例題2）。


\begin{phyomi}[結合エネルギーの曲線]
曲線の高い所ほど、核子1個がしっかりつながっていて、こわれにくい原子核です。いちばん高いのは鉄56・ニッケル62のあたり（約 \qty{8.8}{MeV}。ニッケル62 がわずかに大きい）で、軽い原子核と重い原子核は、それより低くなっています。ヘリウム4（\qty{7.07}{MeV}）は、周りの軽い原子核よりとびぬけて高く、とても安定な原子核です。
\end{phyomi}
\begin{bkfigure}
\begin{tikzpicture}
\begin{axis}[bkplot, width=118mm, height=62mm, xmin=0, xmax=250, ymin=0, ymax=10, xtick={50,100,150,200,250}, ytick={2,4,6,8,10}, xlabel={$A$}, ylabel={$B/A$（MeV）}]
  \fill[bkTheme!10] (axis cs:50,0.05) rectangle (axis cs:66,9.9);
  \fill[bkBlue] (axis cs:2,1.112) circle (1.3pt); \fill[bkBlue] (axis cs:3,2.573) circle (1.3pt); \fill[bkBlue] (axis cs:4,7.074) circle (1.3pt); \fill[bkBlue] (axis cs:6,5.332) circle (1.3pt); \fill[bkBlue] (axis cs:7,5.606) circle (1.3pt); \fill[bkBlue] (axis cs:9,6.463) circle (1.3pt); \fill[bkBlue] (axis cs:10,6.475) circle (1.3pt); \fill[bkBlue] (axis cs:11,6.928) circle (1.3pt); \fill[bkBlue] (axis cs:12,7.680) circle (1.3pt); \fill[bkBlue] (axis cs:14,7.476) circle (1.3pt); \fill[bkBlue] (axis cs:16,7.976) circle (1.3pt); \fill[bkBlue] (axis cs:19,7.779) circle (1.3pt); \fill[bkBlue] (axis cs:20,8.032) circle (1.3pt); \fill[bkBlue] (axis cs:23,8.111) circle (1.3pt); \fill[bkBlue] (axis cs:24,8.261) circle (1.3pt); \fill[bkBlue] (axis cs:27,8.332) circle (1.3pt); \fill[bkBlue] (axis cs:28,8.448) circle (1.3pt); \fill[bkBlue] (axis cs:31,8.481) circle (1.3pt); \fill[bkBlue] (axis cs:32,8.493) circle (1.3pt); \fill[bkBlue] (axis cs:35,8.520) circle (1.3pt); \fill[bkBlue] (axis cs:40,8.595) circle (1.3pt); \fill[bkBlue] (axis cs:40,8.551) circle (1.3pt); \fill[bkBlue] (axis cs:56,8.790) circle (1.3pt); \fill[bkBlue] (axis cs:62,8.795) circle (1.3pt); \fill[bkBlue] (axis cs:63,8.752) circle (1.3pt); \fill[bkBlue] (axis cs:64,8.736) circle (1.3pt); \fill[bkBlue] (axis cs:84,8.717) circle (1.3pt); \fill[bkBlue] (axis cs:88,8.733) circle (1.3pt); \fill[bkBlue] (axis cs:90,8.710) circle (1.3pt); \fill[bkBlue] (axis cs:98,8.635) circle (1.3pt); \fill[bkBlue] (axis cs:107,8.554) circle (1.3pt); \fill[bkBlue] (axis cs:120,8.504) circle (1.3pt); \fill[bkBlue] (axis cs:127,8.445) circle (1.3pt); \fill[bkBlue] (axis cs:132,8.428) circle (1.3pt); \fill[bkBlue] (axis cs:138,8.393) circle (1.3pt); \fill[bkBlue] (axis cs:140,8.376) circle (1.3pt); \fill[bkBlue] (axis cs:142,8.346) circle (1.3pt); \fill[bkBlue] (axis cs:158,8.202) circle (1.3pt); \fill[bkBlue] (axis cs:184,8.005) circle (1.3pt); \fill[bkBlue] (axis cs:197,7.916) circle (1.3pt); \fill[bkBlue] (axis cs:208,7.867) circle (1.3pt); \fill[bkBlue] (axis cs:209,7.848) circle (1.3pt); \fill[bkBlue] (axis cs:232,7.615) circle (1.3pt); \fill[bkBlue] (axis cs:238,7.570) circle (1.3pt);
  \draw[bkGray, line width=0.4pt] (axis cs:5,7.3) -- (axis cs:12,9.1);
  \node[font=\sffamily\footnotesize, anchor=south west, inner sep=1pt] at (axis cs:12,9.0) {${}^4$He};
  \node[font=\sffamily\footnotesize, anchor=west] at (axis cs:7,1.1) {${}^2$H};
  \draw[bkGray, line width=0.4pt] (axis cs:17,7.8) -- (axis cs:27,7.1);
  \node[font=\sffamily\footnotesize, anchor=west, inner sep=1pt] at (axis cs:27,7.0) {${}^{16}$O};
  \node[font=\sffamily\footnotesize, anchor=south] at (axis cs:58,9.05) {${}^{56}$Fe・${}^{62}$Ni};
  \node[font=\sffamily\footnotesize, anchor=north] at (axis cs:236,7.35) {${}^{238}$U};
  \draw[bkRed!75!black, line width=1.1pt, -{Stealth[length=2.4mm]}] (axis cs:15,4.0) -- (axis cs:45,6.0);
  \node[font=\sffamily\footnotesize, text=bkRed!75!black, anchor=north west, align=left] at (axis cs:22,4.0) {核融合\\（軽い核が合わさる）};
  \draw[bkRed!75!black, line width=1.1pt, -{Stealth[length=2.4mm]}] (axis cs:225,6.3) -- (axis cs:130,6.3);
  \node[font=\sffamily\footnotesize, text=bkRed!75!black, anchor=north, align=center] at (axis cs:178,6.1) {核分裂（重い核が2つに分かれる）};
\end{axis}
\end{tikzpicture}
\bkfigcaption{核子1個あたりの結合エネルギー（点は NIST の原子の質量から計算した値）。鉄・ニッケルのあたりがいちばん高く、約 8.8 MeV。矢印の向きに変わると、エネルギーが出る。}\label{ph:原-11:fig-curve}
\end{bkfigure}

この曲線から、原子核からエネルギーを取り出す2つの道がわかります。
\begin{enumerate}[label={\bkcircled{\arabic*}}, leftmargin=2.4\zw, itemsep=1pt, topsep=2pt]
\item \textbf{核分裂}：ウランのような重い原子核（$B/A\approx\qty{7.6}{MeV}$）が、質量数が半分くらいの2つの原子核（$B/A\approx\qty{8.5}{MeV}$）に分かれると、核子1個あたり約 \qty{0.9}{MeV}、236個の核子で約 \qty{200}{MeV} のエネルギーが出ます（原-15）。
\item \textbf{核融合}：水素のような軽い原子核が合わさって、ヘリウムのような曲線の高い原子核になると、エネルギーが出ます。太陽は、これで光っています（⑪）。
\end{enumerate}
どちらも、曲線を\textbf{上へ登る}向きの変化です。鉄より重くする融合や、鉄を分ける分裂では、逆にエネルギーを加えなければなりません。

\Needspace{12\baselineskip}\bksubsection{6}{核反応の式と Q 値}
\label{ph:原-11:sec6}
原子核どうしがぶつかったり、原子核がこわれたりして、違う原子核ができる変化を\textbf{核反応}といいます。\bkidx{かくはんのう}{核反応}核反応は、化学の反応式と同じように、矢印の左に反応の前、右に後を書いて表します。核反応では、次の2つが必ず保存されます。
\begin{enumerate}[label={\bkcircled{\arabic*}}, leftmargin=2.4\zw, itemsep=1pt, topsep=2pt]
\item \textbf{電荷}：左下の数（原子番号。電子は $-1$、陽電子は $+1$）の和が、左右で等しい。
\item \textbf{核子の数}：左上の数（質量数）の和が、左右で等しい。
\end{enumerate}
③のチャドウィックの反応 ${}^{9}_{4}\mathrm{Be}+{}^{4}_{2}\mathrm{He}\to{}^{12}_{6}\mathrm{C}+{}^{1}_{0}\mathrm{n}$ では、左下が $4+2=6+0$、左上が $9+4=12+1$ で、どちらも合っています。

\noindent 1919年、ラザフォードは、窒素に α粒子を当てると陽子が飛び出すことを見つけました。人が初めて元素を変えた反応です。
\[ {}^{14}_{7}\mathrm{N}+{}^{4}_{2}\mathrm{He}\longrightarrow{}^{17}_{8}\mathrm{O}+{}^{1}_{1}\mathrm{p} \]
1932年、イギリスのコッククロフトとウォルトンは、電圧で加速した陽子をリチウムに当てて、ヘリウムの原子核2個に分けました（加速器で原子核を変えた初めての実験。1951年のノーベル物理学賞）。
\[ {}^{7}_{3}\mathrm{Li}+{}^{1}_{1}\mathrm{p}\longrightarrow{}^{4}_{2}\mathrm{He}+{}^{4}_{2}\mathrm{He} \]
飛び出したヘリウムの運動エネルギーの和は、質量の差から $E=mc^2$ で計算した値とよく合いました。$E=mc^2$ が原子核で確かめられたのです。

\phboxhead{Q 値}
核反応で出るエネルギーを \textbf{Q 値}といいます。\bkidx{きゅーち}{Q 値（核反応）}
\pheq{Q=\bigl(\text{反応の前の質量の和}-\text{反応の後の質量の和}\bigr)c^2}{J ＝ kg × (m/s)$^2$}
$Q>0$ ならエネルギーが出て（反応でできた粒の運動エネルギーや γ線になる）、$Q<0$ なら、ぶつける粒の運動エネルギーで足りないぶんを補わなければ、反応は起きません。リチウムの反応では
\[ Q=(7.016003+1.007825-2\times4.002603)\times931.5=\qty{17.3}{MeV} \]
です（陽子の代わりに水素原子の質量を使うと、左右の電子の数が3個ずつでそろいます）。

\begin{phflow}[保存則から、核反応のエネルギーへ]
\phflowrow{エネルギーの保存\\$E=mc^2$}{反応の前と後で\\静止エネルギーを比べる}{$Q=(m_{\text{前}}-m_{\text{後}})c^2$}
\phflowrow{$Q=(m_{\text{前}}-m_{\text{後}})c^2$}{原子の質量で\\電子の数をそろえる}{$Q=931.5\times$（質量の差）\\（$Q$ は MeV、質量は u）}
\end{phflow}

\begin{motto}[3]
\textbf{原子がこわれるときの電子の勘定}\par
中性子は、ひとりでに陽子と電子と反ニュートリノ（⑪）に変わります（β崩壊、原-12）。ニュートリノの質量はとても小さいので、$Q=(m_{\mathrm n}-m_{\mathrm p}-m_{\mathrm e})c^2=\qty{0.782}{MeV}$ です。陽子と電子の質量の和は、水素原子の質量とほぼ同じなので、$Q=(m_{\mathrm n}-m_{\mathrm H})c^2$ と計算しても同じ値になります。炭素14 がβ線を出して窒素14 に変わるとき、原子核の陽子は1個増えますが、出ていく電子も1個なので、\textbf{原子どうしの質量の差}が、そのまま Q 値になります（$(14.003242-14.003074)\times931.5=\qty{0.156}{MeV}$）。
\end{motto}

\begin{phmodel}[原子の質量で、原子核のエネルギーを数える]
\phignore{原子の中の電子をつなぎとめるエネルギー（eV〜数十 keV）。}
\phassume{前と後で電子の数をそろえれば、原子の質量の差が原子核の質量の差になる。}
\phlimit{Q 値が数 keV ほどしかない反応。陽電子を出す反応（電子2個ぶんの \qty{1.022}{MeV} を引く）。}
\end{phmodel}

\phkokomade

\Needspace{12\baselineskip}\phstage{shiki}{光の時計から、相対論を組み立てる}{5}

\bksubsection{7}{光の時計と時間のおくれ、ローレンツ変換}
\label{ph:原-11:sec7}
2枚の鏡を、たてに距離 $L$ はなして向かい合わせ、その間で光を往復させる時計を考えます。光が1往復する時間を1刻みとします。時計といっしょに動く人には、光はまっすぐ上下に動くので、1刻みは
\[ \Delta t_0=\frac{2L}{c} \]
です。この時計が、横に速さ $v$ で動いているのを、地面の人が見るとどうなるでしょう（図\ref{ph:原-11:fig-clock}）。
\begin{bkfigure}
\begin{tikzpicture}[x=1mm, y=0.78mm]
  \begin{scope}
    \txlmirror{(12,4)}{14} \txlmirror{(12,30)}{14}
    \draw[txl light] (11,5) -- (11,29); \draw[txl light] (13,29) -- (13,5);
    \draw[txl dim] (2,4.6) -- (2,29.4); \node[txl lab, anchor=east] at (1,17) {$L$};
    \node[txl lab, anchor=north, align=center] at (12,1.5) {いっしょに動く人\\1刻み $\Delta t_0=2L/c$};
  \end{scope}
  \begin{scope}[shift={(40,0)}]
    \foreach \x in {0,36,72} {\txlmirror{(\x,4)}{12} \txlmirror{(\x,30)}{12}}
    \draw[txl light] (1,5) -- (35,29);
    \draw[txl light] (37,29) -- (71,5);
    \draw[bkInk, txl arr] (78,17) -- (86,17); \node[txl lab, anchor=south] at (82,18) {$v$};
    \draw[txl dim] (0,-1) -- (36,-1); \node[txl lab, anchor=north] at (18,-2) {$v\,\Delta t/2$};
    \node[txl lab, anchor=south east] at (16,18) {$c\,\Delta t/2$};
    \draw[bkGray, dashed, line width=0.5pt] (36,5) -- (36,29);
    \node[txl lab, anchor=west] at (37,12) {$L$};
    \node[txl lab, anchor=north, align=center] at (36,-7) {地面の人：光はななめに進む。1刻みは $\Delta t$};
  \end{scope}
\end{tikzpicture}
\bkfigcaption{光の時計。地面から見ると、光は直角三角形の斜辺を進むので、1刻みが長くなる。}\label{ph:原-11:fig-clock}
\end{bkfigure}
地面の人の1刻みを $\Delta t$ とします。光が下の鏡から上の鏡へ行く間（$\Delta t/2$）に、時計は横へ $v\,\Delta t/2$ 進むので、光は直角三角形の斜辺を進みます。光の速さは地面から見ても $c$ なので、斜辺は $c\,\Delta t/2$ です。三平方の定理（4-4）から
\[ \Bigl(\frac{c\,\Delta t}{2}\Bigr)^2=L^2+\Bigl(\frac{v\,\Delta t}{2}\Bigr)^2 \]
$L=c\,\Delta t_0/2$ を入れて整理すると、$(c^2-v^2)\Delta t^2=c^2\Delta t_0^{\,2}$ なので
\pheq{\Delta t=\gamma\,\Delta t_0,\qquad\gamma=\frac{1}{\sqrt{1-v^2/c^2}}=\frac{1}{\sqrt{1-\beta^2}}}{s ＝ s}
\label{ph:原-11:gamma}となります。$\gamma$ を\textbf{ローレンツ因子}といいます。\bkidx{ろーれんついんし}{ローレンツ因子}$v<c$ なら $\gamma>1$ なので、\textbf{動いている時計は、地面の時計より遅く進んで見えます}。これを\textbf{時間のおくれ}といいます。\bkidx{じかんのおくれ}{時間のおくれ}光の時計だけでなく、原子の時計も、生き物の体の時計も、同じようにおくれます（相対性原理：おくれない時計があれば、それと比べて「自分が動いている」とわかってしまう）。

\begin{phyomi}[$\gamma=1/\sqrt{1-\beta^2}$]
「動いている時計が1秒刻む間に、止まった人の時計は $\gamma$ 秒進む」と読みます。$\beta=v/c$ が小さいと $\gamma\approx1+\tfrac12\beta^2$（5-17）で、ほとんど1です。時速 \qty{1000}{km} の飛行機でも $\gamma-1\approx4\times10^{-13}$ にすぎません。
\end{phyomi}

\begin{phbaisuru}{速さ $v$}
\noindent\begin{minipage}[c]{0.45\linewidth}\raggedright
$v=0.5c$ では $\gamma=1.15$ ですが、速さを光の0.87倍にすると $\gamma=2.0$、0.995倍では $\gamma=10$ です。速さを2倍にしても $\gamma$ は2倍になりません。$c$ に近づくほど急に大きくなり、$v=c$ では無限大になります。物は、どんなに加速しても光の速さには届きません（⑧）。
\end{minipage}\hfill
\begin{minipage}[c]{0.52\linewidth}\centering
\begin{tikzpicture}
\begin{axis}[bkplot, width=58mm, height=44mm, xmin=0, xmax=1.08, ymin=0, ymax=8.5, xtick={0.5,0.87,1}, xticklabels={0.5,0.87,1}, ytick={1,2,4,6,8}, xlabel={$\beta$}, ylabel={$\gamma$}]
  \addplot[bkcurve, domain=0:0.993, samples=160] {1/sqrt(1-x^2)};
  \addplot[bkcurve2, dashed, domain=0:0.95] {1+x^2/2};
  \draw[bkGray, dashed, line width=0.5pt] (axis cs:1,0) -- (axis cs:1,8.3);
  \draw[bkGray, line width=0.5pt] (axis cs:0.5,0) -- (axis cs:0.5,1.155);
  \draw[bkGray, line width=0.5pt] (axis cs:0.866,0) -- (axis cs:0.866,2) -- (axis cs:0,2);
  \node[font=\sffamily\footnotesize, anchor=east, text=bkBlue!80!black] at (axis cs:0.86,5.5) {$\gamma$（実線）};
  \node[font=\sffamily\footnotesize, anchor=west, text=bkOrange!80!black] at (axis cs:0.04,3.2) {$1+\beta^2/2$（破線）};
\end{axis}
\end{tikzpicture}
\end{minipage}
\end{phbaisuru}

\phboxhead{長さの縮み}
地上から見て、高い空で生まれたミューオン（電子の約200倍重い素粒子。止まっていれば平均約 \qty{2.2}{\micro s} でこわれる）が、光の速さに近い速さで地上まで届くのは、時間のおくれのためです（例題5）。ミューオンといっしょに動く人から見ると、ミューオンの時計はおくれていないのに、なぜ地上に届くのでしょうか。その人から見ると、地面のほうが速さ $v$ で近づいてきて、\textbf{動いている物の長さは、動く向きに $1/\gamma$ 倍に縮んで見える}からです。
\pheq{L=\frac{L_0}{\gamma}}{m ＝ m}
$L_0$ は、その物といっしょに止まって見たときの長さです。どちらの人から見ても、ミューオンは地上に届きます。説明のしかたが違うだけで、起こることは同じです。

\phboxhead{ローレンツ変換}
地面の人の座標を $(x,\ t)$、$x$ の向きに速さ $v$ で動くロケットの人の座標を $(x',\ t')$ とします（$t=t'=0$ で原点が重なっていたとします）。ふつうの考え（力-1）では $x'=x-vt$、$t'=t$ ですが、光の速さが両方の人で $c$ になるには、これを少し変えなければなりません。動きはどこでも同じ（一様）なので、変換は1次式で、逆向きも同じ形になるはずです。そこで
\[ x'=\gamma(x-vt),\qquad x=\gamma(x'+vt') \]
とおきます（ロケットの人から見ると、地面は $-v$ で動く）。原点から出た光は、地面では $x=ct$、ロケットでは $x'=ct'$ です。これを入れると $ct'=\gamma(c-v)t$、$ct=\gamma(c+v)t'$。2つをかけ合わせて $tt'$ でわると、$c^2=\gamma^2(c^2-v^2)$ となり、光の時計と同じ $\gamma$ が出てきます。2つめの式に1つめを入れて $t'$ について解くと
\pheq{x'=\gamma(x-vt),\qquad t'=\gamma\Bigl(t-\frac{vx}{c^2}\Bigr)}{m ＝ m、s ＝ s}
\label{ph:原-11:lorentz}が得られます。これを\textbf{ローレンツ変換}といいます。\bkidx{ろーれんつへんかん}{ローレンツ変換}$t'$ の式に $x$ が入っていることが大事で、地面で同時の2つの出来事（同じ $t$、違う $x$）は、ロケットでは同時ではありません。

\begin{motto}[4]
\textbf{「同時」は、見る人によって違う}\par
ロケットのまん中で、前と後ろに向けて同時に光を出します。ロケットの人には、光は前と後ろの壁に同時に届きます。ところが、地面の人から見ると、後ろの壁は光に向かって近づき、前の壁は光から遠ざかるので、光は後ろの壁に先に届きます。光の速さはどちらの人にも $c$ なので、どちらの見方も正しいのです。ローレンツ変換の $t'=\gamma(t-vx/c^2)$ に $x$ が入っているのは、このためです。ただし、一方の出来事が原因でもう一方が起きる（光の速さ以下で知らせが届く）2つの出来事では、順番が入れかわることはありません。
\end{motto}

\phboxhead{速度の合成}
ロケットの中で、物が前へ速さ $u'$ で動きます（$x'=u't'$）。地面から見た速さ $u=x/t$ は、逆向きの変換 $x=\gamma(x'+vt')$、$t=\gamma(t'+vx'/c^2)$ に $x'=u't'$ を入れて
\pheq{u=\frac{x}{t}=\frac{u'+v}{1+u'v/c^2}}{m/s ＝ m/s}
\label{ph:原-11:gousei}です。$u'=c$ を入れると $u=(c+v)/(1+v/c)=c$ となり、光の速さは、やはり $c$ です。

\Needspace{12\baselineskip}\begin{phflow}[ローレンツ変換から、力学の式へ]
\phflowrow{$x'=\gamma(x-vt)$\\$t'=\gamma(t-vx/c^2)$}{$v\ll c$ なので\\$\gamma\approx1$、$vx/c^2\approx0$}{$x'=x-vt$、$t'=t$}
\phflowrow{$u=\dfrac{u'+v}{1+u'v/c^2}$}{$u'v\ll c^2$}{$u=u'+v$（力-1）}
\end{phflow}

\begin{phmijika}[GPS の衛星の時計]
GPS の衛星は、地球の中心から約 \qty{26560}{km}（高さ約 \qty{20200}{km}、波-14）を、秒速約 \qty{3.9}{km} で回っています。この速さで、衛星の時計は地上より1日に約 \qty{7}{\micro s} おくれます（例題6）。いっぽう、重力が弱い高い所の時計は速く進むことが、一般相対性理論からわかっています（1日に約 \qty{46}{\micro s}）。合わせて約 \qty{38}{\micro s} 進むので、GPS は衛星の時計を前もって遅らせています。\qty{38}{\micro s} の間に、電波は 11〜12 km ほど進みます。補正しなければ、位置が1日に何 km もずれてしまいます。
\end{phmijika}

\Needspace{12\baselineskip}\bksubsection{8}{相対論の運動量とエネルギー}
\label{ph:原-11:sec8}
光の速さに近い物では、運動量とエネルギーも、ニュートンの力学の形から変わります。どの慣性系でも「ぶつかる前と後で運動量の和が変わらない」ことが成り立つように決めると、質量 $m$・速さ $v$ の物の運動量とエネルギーは
\pheq{p=\gamma mv,\qquad E=\gamma mc^2}{kg·m/s、J}
\label{ph:原-11:pE}になります（力-2・力-3 で「光の速さに近い物の運動量は $mv$ ではなくなる」と書いたのは、このことです）。$\gamma$ は $v\to c$ で無限大になるので、速さが $c$ に近づくと、運動量もエネルギーも、どこまでも大きくなります。だから、有限のエネルギーでは、物を光の速さまで加速できません。

2つの式から $\gamma$ と $v$ を消してみます。
\[ E^2-(pc)^2=\gamma^2m^2c^4-\gamma^2m^2v^2c^2=\gamma^2m^2c^4\Bigl(1-\frac{v^2}{c^2}\Bigr)=m^2c^4 \]
$\gamma^2(1-v^2/c^2)=1$ を使いました。こうして、速さによらない関係
\pheq{E^2=(pc)^2+(mc^2)^2}{J$^2$ ＝ J$^2$ ＋ J$^2$}
\label{ph:原-11:Ep}が得られます。直角三角形の3辺のような形で、斜辺が $E$、2つの辺が $pc$ と $mc^2$ です。

\begin{phflow}[$E^2=(pc)^2+(mc^2)^2$ から]
\phflowrow{$E^2=(pc)^2+(mc^2)^2$}{止まっている\\（$p=0$）}{$E=mc^2$}
\phflowrow{$E=\gamma mc^2$}{$v\ll c$ なので\\$\gamma\approx1+\tfrac12\beta^2$}{$E\approx mc^2+\tfrac12mv^2$}
\phflowrow{$E^2=(pc)^2+(mc^2)^2$}{質量が0\\（光子）}{$E=pc$、$p=\dfrac{E}{c}=\dfrac{h}{\lambda}$}
\end{phflow}
\label{ph:原-11:emc2}1段めが②の $E=mc^2$ です。2段めは、エネルギーが「静止エネルギー＋ニュートンの運動エネルギー」になることを示しています。だから、運動エネルギーは
\pheq{K=E-mc^2=(\gamma-1)mc^2}{J ＝ J}
と決めます。3段めは、質量のない光子では $E=pc$ になることを示しています。これは、原-2 の光子の運動量 $p=h/\lambda$ や、電-15 の光の運動量 $p=U/c$ と同じです。

\begin{phbaisuru}{運動エネルギー $K$}
\noindent\begin{minipage}[c]{0.45\linewidth}\raggedright
電子の運動エネルギーを増やしていくと、ニュートンの式 $K=\tfrac12mv^2$（破線）では、速さはどこまでも大きくなり、$K=\qty{0.26}{MeV}$ で光の速さをこえてしまいます。相対論の式（実線）では、$K=mc^2=\qty{0.511}{MeV}$ で $v=0.87c$、$K$ を2倍にしても $0.94c$ で、光の速さに近づくだけです。
\end{minipage}\hfill
\begin{minipage}[c]{0.52\linewidth}\centering
\begin{tikzpicture}
\begin{axis}[bkplot, width=58mm, height=44mm, xmin=0, xmax=2.2, ymin=0, ymax=1.65, xtick={0.5,1,1.5,2}, ytick={0.5,1,1.5}, xlabel={$K$（MeV）}, ylabel={$v/c$}]
  \addplot[bkcurve, domain=0:2.1, samples=120] {sqrt(1-1/(1+x/0.511)^2)};
  \addplot[bkcurve2, dashed, domain=0:0.6, samples=60] {sqrt(2*x/0.511)};
  \draw[bkGray, dashed, line width=0.5pt] (axis cs:0,1) -- (axis cs:2.1,1);
  \node[font=\sffamily\footnotesize, anchor=north, text=bkBlue!80!black] at (axis cs:1.5,0.9) {相対論（実線）};
  \node[font=\sffamily\footnotesize, anchor=west, text=bkOrange!80!black] at (axis cs:0.62,1.45) {ニュートン（破線）};
\end{axis}
\end{tikzpicture}
\end{minipage}
\end{phbaisuru}

\begin{bkfigure}
\begin{tikzpicture}
\begin{axis}[bkplot, width=86mm, height=52mm, xmin=0, xmax=3.2, ymin=0, ymax=3.4, xtick={1,2,3}, ytick={0.511,1,2,3}, yticklabels={0.511,1,2,3}, xlabel={$pc$（MeV）}, ylabel={$E$（MeV）}]
  \addplot[bkcurve, domain=0:3.0, samples=80] {sqrt(x^2+0.511^2)};
  \addplot[bkcurve2, dashed, domain=0:3.0] {x};
  \draw[bkGray, line width=0.5pt] (axis cs:1.422,0) -- (axis cs:1.422,1.511) -- (axis cs:0,1.511);
  \node[font=\sffamily\footnotesize, anchor=south east, text=bkBlue!80!black] at (axis cs:1.35,1.6) {電子（実線）};
  \node[font=\sffamily\footnotesize, anchor=north west, text=bkOrange!80!black] at (axis cs:2.0,1.9) {光子 $E=pc$（破線）};
  \node[font=\sffamily\footnotesize, anchor=north west] at (axis cs:1.45,1.0) {例題9};
\end{axis}
\end{tikzpicture}
\bkfigcaption{電子のエネルギー $E$ と運動量 $pc$。$p=0$ では $E=mc^2=\qty{0.511}{MeV}$。運動量が大きくなると、光子の直線 $E=pc$ に近づく。}\label{ph:原-11:fig-Ep}
\end{bkfigure}

\phboxhead{コンプトンの式を導く}
原-2 では、X線が電子にぶつかって波長がのびるコンプトン散乱を、式の結果だけ読みました。ここで、エネルギーと運動量の保存から導きます（図\ref{ph:原-11:fig-compton}）。
\begin{bkfigure}
\begin{tikzpicture}[x=1mm, y=1mm]
  \draw[txl light] (0,12) -- (30,12); \node[txl lab, anchor=south] at (15,13) {光子 $E$、$p=E/c$};
  \fill[bkBlue!70] (34,12) circle (1.6); \node[txl lab, anchor=north east] at (33,10) {止まった電子};
  \draw[txl light] (37,13.5) -- (62,27); \node[txl lab, anchor=south east] at (52,23) {光子 $E'$};
  \draw[bkBlue!80!black, line width=1.1pt, -{Stealth[length=2.2mm]}] (37,10.5) -- (60,0);
  \node[txl lab, anchor=west] at (61,0) {電子 $E_{\mathrm e}$、$\vec p_{\mathrm e}$};
  \draw[bkGray, dashed, line width=0.5pt] (37,12) -- (66,12);
  \draw[bkInk, line width=0.5pt] (47,12) arc (0:28:10); \node[txl lab, anchor=west] at (47.5,15.5) {$\theta$};
\end{tikzpicture}
\bkfigcaption{コンプトン散乱。エネルギー $E$ の光子が、止まった電子にぶつかり、角 $\theta$ の向きへエネルギー $E'$ で散らばる。}\label{ph:原-11:fig-compton}
\end{bkfigure}
運動量の保存は、ベクトルで $\vec p=\vec p\,'+\vec p_{\mathrm e}$ です。$\vec p_{\mathrm e}=\vec p-\vec p\,'$ の大きさの2乗を余弦定理（4-7）で書くと
\[ (p_{\mathrm e}c)^2=E^2+E'^2-2EE'\cos\theta \]
エネルギーの保存は $E+mc^2=E'+E_{\mathrm e}$ なので、$E_{\mathrm e}=E-E'+mc^2$ を2乗して、$E_{\mathrm e}^{\,2}=(p_{\mathrm e}c)^2+(mc^2)^2$ に入れます。
\[ (E-E')^2+2(E-E')mc^2+(mc^2)^2=E^2+E'^2-2EE'\cos\theta+(mc^2)^2 \]
左右の $E^2$・$E'^2$・$(mc^2)^2$ を消して整理すると、$(E-E')mc^2=EE'(1-\cos\theta)$ です。両辺を $EE'mc^2$ でわり、$E=hc/\lambda$ を使うと
\pheq{\lambda'-\lambda=\frac{h}{mc}(1-\cos\theta)}{m ＝ J·s/(kg·m/s)}
\label{ph:原-11:compton}が出てきます。これが\textbf{コンプトンの式}です。\bkidx{こんぷとんのしき}{コンプトンの式}$h/(mc)=\qty{2.43}{pm}$ を電子のコンプトン波長といいます。光子を「エネルギー $E$・運動量 $E/c$ を持つ粒」と考え、電子を相対論で扱うと、実験とぴったり合うのです（例題10）。力-3 の運動量の保存は、光の速さに近い粒の衝突でも、そのまま使えます。

\Needspace{8\baselineskip}\bksubsection{9}{光のドップラー効果}
\label{ph:原-11:sec9}
波-8 では、音のドップラー効果を学び、光では近似の式 $\Delta f\approx vf/c$ だけを使いました。光には、音にとっての空気のような「波を伝える物」がないので、正しい式は相対性理論で求めます。

光源が、見る人から速さ $v$ で遠ざかりながら、自分の時計で周期 $T_0=1/f_0$ ごとに波の山を出すとします。見る人の時計では、時間のおくれで、山と山の間は $\gamma T_0$ です。その間に光源は $v\gamma T_0$ だけ遠ざかるので、次の山は、その距離を進むぶん $v\gamma T_0/c$ だけおくれて届きます。届く山の間かくは
\[ T=\gamma T_0(1+\beta)=T_0\frac{1+\beta}{\sqrt{(1-\beta)(1+\beta)}}=T_0\sqrt{\frac{1+\beta}{1-\beta}} \]
なので、見る人の受けとる振動数は
\pheq{f=f_0\sqrt{\frac{1-\beta}{1+\beta}}\qquad\text{（遠ざかるとき。近づくときは $\beta$ を $-\beta$ に）}}{Hz ＝ Hz}
\label{ph:原-11:doppler}です。遠ざかる星の光は振動数が小さく（波長が長く、赤いほうに）なり、近づく星の光は振動数が大きく（青いほうに）なります。

\begin{phflow}[波-8 の近似へ]
\phflowrow{$f=f_0\sqrt{\dfrac{1-\beta}{1+\beta}}$}{$\beta\ll1$ なので\\$\beta^2$ を無視する}{$f\approx f_0(1-\beta)$、$\Delta f\approx\dfrac{v}{c}f_0$}
\end{phflow}

\begin{phbaisuru}{遠ざかる速さ $\beta$}
\noindent\begin{minipage}[c]{0.44\linewidth}\raggedright
$\beta=0.1$ では、相対論の式 0.905、音の式（見る人が止まり、音源が遠ざかる形）0.909、近似 0.900 で、ほとんど同じです。$\beta=0.6$ では、相対論 0.50、音の式 0.625、近似 0.40 と大きく違ってきます。速さを2倍にしても、ずれは2倍よりも大きくなります。
\end{minipage}\hfill
\begin{minipage}[c]{0.53\linewidth}\centering
\begin{tikzpicture}
\begin{axis}[bkplot, width=58mm, height=44mm, xmin=0, xmax=0.88, ymin=0, ymax=1.12, xtick={0.2,0.4,0.6,0.8}, ytick={0.5,1}, xlabel={$\beta$}, ylabel={$f/f_0$}]
  \addplot[bkcurve, domain=0:0.85] {sqrt((1-x)/(1+x))};
  \addplot[bkcurve2, dashed, domain=0:0.85] {1/(1+x)};
  \addplot[bkcurve3, densely dotted, domain=0:0.85] {1-x};
  \node[font=\sffamily\footnotesize, anchor=west, text=bkBlue!80!black] at (axis cs:0.03,0.42) {相対論（実線）};
  \node[font=\sffamily\footnotesize, anchor=west, text=bkOrange!80!black] at (axis cs:0.03,0.28) {音の式（破線）};
  \node[font=\sffamily\footnotesize, anchor=west, text=bkGreen!50!black] at (axis cs:0.03,0.14) {近似（点線）};
\end{axis}
\end{tikzpicture}
\end{minipage}
\end{phbaisuru}

\begin{phmijika}[星の動きをはかる]
水素の原子は、決まった波長 \qty{656.3}{nm} の赤い光を出します（原-4）。遠くの星の光を分光器（波-13）で分けて、この線の波長がどれだけずれているかを調べると、星がこちらに近づいているか、遠ざかっているか、どれくらいの速さかがわかります。星の周りを回る惑星に引かれて、星がわずかにゆれることも、この方法で見つけられます。
\end{phmijika}

\Needspace{12\baselineskip}\bksubsection{10}{核力と半経験的質量公式}
\label{ph:原-11:sec10}
原子核の中では、陽子どうしが約 \qty{1}{fm} まで近づいています。そのクーロン力は約 \qty{230}{N} の反発力です（電-1 の章末問題、例題12）。陽子1個の重さの $10^{28}$ 倍もの力で押し合っているのに、原子核はばらばらになりません。陽子と中性子を引きつけて離さない、別の強い力がはたらいているのです。これを\textbf{核力}といいます。\bkidx{かくりょく}{核力}

\noindent\begin{minipage}[c]{0.5\linewidth}\raggedright
核力には、次の特徴があります。
\begin{enumerate}[label={\bkcircled{\arabic*}}, leftmargin=2.4\zw, itemsep=1pt, topsep=2pt]
\item 約 1〜2 fm の近い所では、クーロン力よりずっと強く引き合う。
\item それより遠いと、急に0になる（届く距離が短い）。
\item 陽子と陽子、陽子と中性子、中性子と中性子のどれにも、ほぼ同じようにはたらく。
\item 近づきすぎると、逆に強くはね返す。
\end{enumerate}
\end{minipage}\hfill
\begin{minipage}[c]{0.47\linewidth}\centering
\begin{tikzpicture}
\begin{axis}[bkplot, axis lines=left, width=50mm, height=44mm, xmin=0, xmax=4.2, ymin=-42, ymax=45, xtick={1,2,3,4}, ytick={0}, xlabel={$r$（fm）}, ylabel={位置エネルギー}]
  \draw[bkGray, line width=0.5pt] (axis cs:0,0) -- (axis cs:4.2,0);
  \addplot[bkcurve2, dashed, domain=0.6:4.0, samples=80] {1.44/x*8};
  \addplot[bkcurve, domain=0.72:4.0, samples=160] {300*exp(-(x-0.5)/0.14)-80*exp(-(x-0.5)/0.75)+1.44/x*8};
  \node[font=\sffamily\footnotesize, anchor=west, text=bkOrange!80!black] at (axis cs:2.0,22) {クーロンだけ（破線）};
  \node[font=\sffamily\footnotesize, anchor=west, text=bkBlue!80!black] at (axis cs:1.9,-19) {核力も入れる（実線）};
\end{axis}
\end{tikzpicture}
\end{minipage}
\par\vspace{3pt}
1935年、湯川秀樹は、核力は核子の間で、電子の約200倍の質量を持つ新しい粒（中間子）がやりとりされて生まれる、と考えました。力を運ぶ粒の質量が $m$ なら、力の届く距離はおよそ $\hbar/(mc)$ になります（$\hbar=h/(2\pi)$）。のちに見つかったπ中間子の静止エネルギーは約 \qty{140}{MeV} で、届く距離は
\[ \frac{\hbar}{mc}=\frac{\hbar c}{mc^2}=\frac{\qty{197.3}{MeV.fm}}{\qty{140}{MeV}}\approx\qty{1.4}{fm} \]
と、核力の届く距離に合います。湯川は1949年に、日本人で初めてノーベル物理学賞を受けました。

\phboxhead{液滴模型と半経験的質量公式}
核力は近くの核子どうしにしかはたらかず、原子核の密度はどれもほぼ同じでした（③）。これは、分子が近くの分子とだけ引き合う水の\textbf{しずく}とよく似ています。原子核を、電気を帯びたしずくとみて、結合エネルギーを5つの項で書いた式が\textbf{半経験的質量公式}です。\bkidx{はんけいけんてきしつりょうこうしき}{半経験的質量公式}
\pheq{B=a_{\mathrm V}A-a_{\mathrm S}A^{2/3}-a_{\mathrm C}\frac{Z(Z-1)}{A^{1/3}}-a_{\mathrm A}\frac{(A-2Z)^2}{A}+\delta}{MeV}
\label{ph:原-11:semf}
\begin{enumerate}[label={\bkcircled{\arabic*}}, leftmargin=2.4\zw, itemsep=1pt, topsep=2pt]
\item \textbf{体積の項} $a_{\mathrm V}A$：核子は近くの核子とだけ結びつくので、結合エネルギーは核子の数 $A$ に比例する。
\item \textbf{表面の項} $-a_{\mathrm S}A^{2/3}$：表面の核子は、外側に相手がいないぶん結びつきが弱い。表面積は $R^2\propto A^{2/3}$。
\item \textbf{クーロンの項}：陽子 $Z$ 個のどの2個も押し合う。組の数は $Z(Z-1)/2$、距離は $R\propto A^{1/3}$。
\item \textbf{対称の項}：陽子と中性子の数がそろっているほど安定（原-7 のパウリの排他律による）。
\item \textbf{対の項} $\delta$：陽子・中性子がどちらも偶数なら $+$、どちらも奇数なら $-$。
\end{enumerate}
係数は、たくさんの原子核の質量に合うように決めた値で、教科書によって少しずつ違います。ここでは一例として $a_{\mathrm V}=15.8$、$a_{\mathrm S}=18.3$、$a_{\mathrm C}=0.714$、$a_{\mathrm A}=23.2$（いずれも MeV）、$\delta=\pm12.0/\sqrt A$ MeV を使います。

\begin{bkfigure}
\begin{tikzpicture}
\begin{axis}[bkplot, width=118mm, height=58mm, xmin=0, xmax=250, ymin=0, ymax=17, xtick={50,100,150,200,250}, ytick={4,8,12,16}, xlabel={$A$}, ylabel={$B/A$（MeV）}]
  \addplot[bkthin, domain=4:245] {15.8};
  \addplot[bkcurve2, dashed, domain=8:245, samples=80] {15.8-18.3/x^(1/3)};
  \addplot[bkcurve3, densely dotted, domain=8:245, samples=80] {15.8-18.3/x^(1/3)-0.714*(x/(2+0.0153*x^(2/3)))*(x/(2+0.0153*x^(2/3))-1)/x^(4/3)};
  \addplot[bkcurve, domain=12:240, samples=80] coordinates {(12,7.316) (16,7.733) (20,8.009) (24,8.196) (28,8.326) (32,8.461) (36,8.585) (40,8.667) (44,8.719) (48,8.749) (52,8.761) (56,8.761) (60,8.763) (64,8.777) (68,8.779) (72,8.770) (76,8.752) (80,8.728) (84,8.718) (88,8.704) (92,8.683) (96,8.657) (100,8.627) (104,8.609) (108,8.585) (112,8.557) (116,8.524) (120,8.499) (124,8.473) (128,8.443) (132,8.409) (136,8.382) (140,8.354) (144,8.322) (148,8.288) (152,8.261) (156,8.231) (160,8.198) (164,8.167) (168,8.138) (172,8.106) (176,8.073) (180,8.044) (184,8.014) (188,7.981) (192,7.951) (196,7.921) (200,7.889) (204,7.859) (208,7.829) (212,7.797) (216,7.767) (220,7.738) (224,7.707) (228,7.677) (232,7.648) (236,7.617) (240,7.588)};
  \fill[bkInk] (axis cs:4,7.074) circle (1.1pt); \fill[bkInk] (axis cs:12,7.680) circle (1.1pt); \fill[bkInk] (axis cs:16,7.976) circle (1.1pt); \fill[bkInk] (axis cs:20,8.032) circle (1.1pt); \fill[bkInk] (axis cs:28,8.448) circle (1.1pt); \fill[bkInk] (axis cs:40,8.551) circle (1.1pt); \fill[bkInk] (axis cs:56,8.790) circle (1.1pt); \fill[bkInk] (axis cs:62,8.795) circle (1.1pt); \fill[bkInk] (axis cs:90,8.710) circle (1.1pt); \fill[bkInk] (axis cs:120,8.504) circle (1.1pt); \fill[bkInk] (axis cs:138,8.393) circle (1.1pt); \fill[bkInk] (axis cs:158,8.202) circle (1.1pt); \fill[bkInk] (axis cs:184,8.005) circle (1.1pt); \fill[bkInk] (axis cs:208,7.867) circle (1.1pt); \fill[bkInk] (axis cs:238,7.570) circle (1.1pt);
  \node[font=\sffamily\footnotesize, anchor=south] at (axis cs:150,15.9) {① 体積だけ（細い実線）};
  \node[font=\sffamily\footnotesize, anchor=south, text=bkOrange!80!black] at (axis cs:165,12.9) {①＋② 表面を引く（破線）};
  \node[font=\sffamily\footnotesize, anchor=south, text=bkGreen!50!black] at (axis cs:175,9.6) {③ クーロンも引く（点線）};
  \node[font=\sffamily\footnotesize, anchor=north, text=bkBlue!80!black] at (axis cs:120,7.7) {④⑤も入れた式（太い実線）・点は NIST};
\end{axis}
\end{tikzpicture}
\bkfigcaption{半経験的質量公式の項を、1つずつ引いていくようす（核子1個あたり。$Z$ は、その $A$ でいちばん安定な値）。軽い所は表面の項、重い所はクーロンの項で下がり、山ができる。}\label{ph:原-11:fig-semf}
\end{bkfigure}
図\ref{ph:原-11:fig-semf}のように、軽い原子核は表面の割合が大きいので $B/A$ が小さく、重い原子核は陽子どうしの反発が大きいので $B/A$ が小さくなります。その間で、鉄のあたりに山ができるのです。ウラン236 が、半分の大きさの原子核2つに分かれると、表面は増えますが、クーロンの反発がずっと減るので、式から約 \qty{180}{MeV} のエネルギーが出ると計算できます（原-15）。

\begin{phcode}[Python]
def B(A, Z):    \# 半経験的質量公式（MeV）
    aV, aS, aC, aA, aP = 15.8, 18.3, 0.714, 23.2, 12.0
    N = A - Z
    d = 0
    if Z \% 2 == 0 and N \% 2 == 0: d = aP / A**0.5
    if Z \% 2 == 1 and N \% 2 == 1: d = -aP / A**0.5
    return (aV*A - aS*A**(2/3) - aC*Z*(Z-1)/A**(1/3)
            - aA*(A-2*Z)**2/A + d)
for A, Z in [(16, 8), (56, 26), (238, 92)]:
    print(A, round(B(A, Z), 1), round(B(A, Z)/A, 2))
\end{phcode}
\begin{phcodeout}[実行結果]
16 123.7 7.73
56 490.6 8.76
238 1809.6 7.6
\end{phcodeout}
左から、質量数、結合エネルギー（MeV）、核子1個あたり（MeV）です。NIST の質量から計算した本当の値は、酸素16 が 127.6、鉄56 が 492.3、ウラン238 が 1801.7 MeV です。5つの項だけで、たくさんの原子核の結合エネルギーを、重い原子核では 1\,\% より小さいずれで表せます。

\begin{phmodel}[液滴模型]
\phignore{原子核の中の核子の一つ一つの動き（殻の構造）。原子核の形のゆがみ。}
\phassume{原子核は、密度が一定の電気を帯びたしずく。結合エネルギーは体積・表面・クーロン・対称・対の5つの項の和。}
\phlimit{軽い原子核。陽子や中性子の数が 2・8・20・28・50・82・126（魔法数）の原子核は、式より強く結びついている（酸素16 で約 \qty{4}{MeV}）。}
\end{phmodel}

\Needspace{12\baselineskip}\bksubsection{11}{太陽の核融合と素粒子}
\label{ph:原-11:sec11}
太陽の中心は、温度が約1500万 K の、水素の原子核（陽子）と電子のガスです。そこでは、陽子どうしが次のようにくっついて、最後にヘリウム4 ができます（\textbf{核融合}。\bkidx{かくゆうごう}{核融合}図\ref{ph:原-11:fig-pp}）。
\begin{bkfigure}
\begin{tikzpicture}[x=1mm, y=1mm]
  % ① 陽子＋陽子 → D ＋ e+ ＋ ν
  \node[txl lab, anchor=east] at (-2,40) {①};
  \txlp{(4,40)}{2.2} \node[txl lab] at (9,40) {＋}; \txlp{(14,40)}{2.2}
  \draw[bkInk, txl arr] (19,40) -- (29,40);
  \txln{(37.4,40)}{2.2} \txlp{(33,40)}{2.2}
  \node[txl lab, anchor=west] at (41,40) {重水素 D ＋ 陽電子 e$^+$ ＋ ニュートリノ $\nu$};
  \node[txl lab, anchor=west] at (104,40) {（2回）};
  % ② D ＋ 陽子 → 3He ＋ γ線
  \node[txl lab, anchor=east] at (-2,24) {②};
  \txln{(8.4,24)}{2.2} \txlp{(4,24)}{2.2} \node[txl lab] at (13.5,24) {＋}; \txlp{(18.5,24)}{2.2}
  \draw[bkInk, txl arr] (23,24) -- (31,24);
  \txln{(37.2,22.1)}{2.2} \txlp{(35,25.9)}{2.2} \txlp{(39.4,25.9)}{2.2}
  \node[txl lab, anchor=west] at (43,24) {ヘリウム3 ＋ γ線};
  \node[txl lab, anchor=west] at (104,24) {（2回）};
  % ③ 3He ＋ 3He → 4He ＋ 陽子2個
  \node[txl lab, anchor=east] at (-2,8) {③};
  \txln{(5.2,6.1)}{2.2} \txlp{(3,9.9)}{2.2} \txlp{(7.4,9.9)}{2.2} \node[txl lab] at (12.5,8) {＋};
  \txln{(20.2,6.1)}{2.2} \txlp{(18,9.9)}{2.2} \txlp{(22.4,9.9)}{2.2}
  \draw[bkInk, txl arr] (27,8) -- (35,8);
  \txln{(43.4,10)}{2.2} \txln{(39,5.6)}{2.2} \txlp{(39,10)}{2.2} \txlp{(43.4,5.6)}{2.2}
  \node[txl lab, anchor=west] at (47,8) {ヘリウム4 ＋};
  \txlp{(73,8)}{2.2} \txlp{(78.5,8)}{2.2}
  \node[txl lab, anchor=west] at (82,8) {（陽子2個）};
\end{tikzpicture}
\bkfigcaption{太陽の中心の核融合（いちばん多い道すじ。＋の印の丸が陽子、印のない丸が中性子）。①と②が2回ずつ起きて、③でヘリウム4 ができる。}\label{ph:原-11:fig-pp}
\end{bkfigure}
全部まとめると、陽子4個がヘリウム4 になり、陽電子2個とニュートリノ2個ができます。陽電子は周りの電子とすぐに出会って消え、γ線になります。

陽子どうしは、近づくとクーロン力で強く押し返されます。核力が届く 1〜2 fm まで近づくには、1 MeV ほどのエネルギーが要りますが、1500万 K の熱の運動のエネルギーの目安 $k_{\mathrm B}T$ は、たった 1.3 keV ほどです。それでも反応が起きるのは、原-6 で学んだトンネル効果で、クーロン力の壁をわずかな確率ですりぬけるからです。この確率がとても小さいので、太陽はゆっくり燃え、何十億年も光り続けられるのです。

\phboxhead{\bkhatenaref の答え}
\label{ph:原-11:hatena}太陽は、中心で水素の原子核（陽子）4個をヘリウム4 の原子核1個に変える\textbf{核融合}で光っています。できたヘリウム4 は、もとの水素4個より 0.71\,\% 軽く、その質量の差が $E=mc^2$ で、1回あたり \qty{26.7}{MeV} のエネルギーになります（例題4）。太陽の明るさ $\qty{3.83e26}{W}$ から、太陽は1秒に約 $\qty{4.3e9}{kg}$ ずつ軽くなり、約 $\qty{6.0e11}{kg}$ の水素をヘリウムに変えていることがわかります。質量あたりのエネルギーが石炭の約2千万倍も大きいので、石炭なら5000年で燃えつきるところを、太陽は100億年ほども光り続けられると考えられています。

\begin{phmijika}[地上の太陽：核融合の実験炉]
地上で核融合の電気をつくる研究では、重水素（D）と三重水素（T）の反応 $\mathrm D+\mathrm T\to{}^4\mathrm{He}+\mathrm n$（\qty{17.6}{MeV}、例題3）を使います。陽子どうしより、はるかに反応しやすいからです。日本もふくむ国々が南フランスで建設している実験炉 ITER は、外から入れる \qty{50}{MW} の10倍の \qty{500}{MW} を、核融合で出すことをめざしています（物理シリーズの『熱学』（熱-16）、原-15）。
\end{phmijika}

\phboxhead{素粒子}
\label{ph:原-11:soryushi}陽子と中性子は、さらに小さな\textbf{クォーク}が3個集まってできています。陽子は u（アップ）クォーク2個と d（ダウン）クォーク1個、中性子は u 1個と d 2個です。電子やニュートリノは、それ以上分けられないと考えられている\textbf{レプトン}です。\bkidx{そりゅうし}{素粒子}このように、物をつくるいちばん小さな粒を\textbf{素粒子}といいます。
\begin{bkcenter}
\begin{bktable}{colspec={Q[l,m,wd=26mm] Q[l,m,wd=50mm] Q[l,m,wd=40mm]}, rowsep=0.5pt, row{1}={bg=bkTheme!10}, cells={font=\small}}
力 & 何にはたらくか & 運ぶ粒・身近な例 \\
重力 & 質量のあるすべて（いちばん弱い） & 星や惑星の運動 \\
電磁気力 & 電荷を持つ粒 & 光子。原子・化学の反応 \\
強い力 & クォーク（核力のもと） & グルーオン。原子核 \\
弱い力 & クォークとレプトン & W・Z 粒子。β崩壊、太陽の最初の反応 \\
\end{bktable}
\bktabcaption{自然の4つの力}\label{ph:原-11:tab-force}
\end{bkcenter}
太陽の最初の反応（陽子2個 → 重水素）では、陽子が中性子に変わります。これは弱い力による反応で、とても起こりにくいことも、太陽がゆっくり燃えるわけの1つです。

どの粒にも、質量が同じで電荷が反対の\textbf{反粒子}があります。電子の反粒子は陽電子です。電子と陽電子が止まって出会うと、2つとも消えて、\qty{0.511}{MeV} の γ線が2本、反対向きに出ます（$2m_{\mathrm e}c^2=\qty{1.022}{MeV}$）。病院の PET 検査は、陽電子を出す薬を体に入れ、この2本の γ線を、体の周りの検出器で同時にとらえて、薬が集まった場所を写します。

太陽の核融合で生まれたニュートリノは、地球もほとんど素通りします。岐阜県の神岡の地下の大きな水そうで、宇宙からのニュートリノをとらえた研究（2002年）と、ニュートリノに質量があることを示した研究（2015年）は、どちらも日本の研究者がノーベル物理学賞を受けました。

\bksubsection{12}{時間と空間を1つに}
\label{ph:原-11:sec12}
ローレンツ変換では、時間 $t$ と位置 $x$ が混ざり合いました。そこで、専門の本では、$(ct,\ x,\ y,\ z)$ の4つを1つのベクトルとみなし、$E^2-(pc)^2$ が「どの慣性系でも同じ長さ」になるように、エネルギーと運動量も $(E/c,\ \vec p)$ の4つの組にまとめます。電-14 で見た電場と磁場の入れかわりも、電場と磁場を1つの表（電磁場テンソル）にまとめると、ローレンツ変換でまとめて変わる形に書けます。下の専門書の書き方を、ちょっとのぞいてみましょう。

\begin{senmonsho}[相対論の専門書の書き方]
\begin{bkhonbun}
4 元運動量を $p^\mu=(E/c,\,\boldsymbol p)$ とすると，そのミンコフスキー計量のノルムはローレンツ不変であり，
\[ p_\mu p^\mu=\frac{E^2}{c^2}-\boldsymbol p^2=m^2c^2 \]
が成り立つ．同様に，電磁場は反対称テンソル $F^{\mu\nu}$ にまとめられ，$F'^{\mu\nu}=\Lambda^\mu{}_\alpha\Lambda^\nu{}_\beta F^{\alpha\beta}$ と変換する．
\end{bkhonbun}
\bkyaku{「4元」は4つの成分。$\mu$ は 0〜3 の番号で、0 が時間、1〜3 が空間。「ノルム」はベクトルの長さの2乗にあたるもの（ミンコフスキー計量では、時間の成分の2乗から空間の成分の2乗を引く）。「ローレンツ不変」は、どの慣性系で計算しても同じ値になること。この章の $E^2=(pc)^2+(mc^2)^2$ を $c^2$ でわった式です。$\Lambda$ はローレンツ変換の表（行列）。$\boldsymbol p$ の太字は、この本の $\vec p$。}
\end{senmonsho}

%% ---------------------------------------------------------------
\begin{phmada}[ニュートリノの質量は、いくら？]
ニュートリノは、長い間、質量が0だと考えられていました。ところが、飛んでいる間に種類が入れかわる「ニュートリノ振動」が見つかり、質量を持つことがわかりました（2015年のノーベル物理学賞）。でも、振動からわかるのは種類どうしの質量の差だけで、質量そのものの値は、まだはかれていません。ドイツの KATRIN 実験は、三重水素のβ崩壊で出る電子のエネルギーを細かく調べ、電子ニュートリノの質量は \qty{0.45}{eV} より小さい（電子の100万分の1より軽い）ことを示しました（2024年の発表）。なぜ、ほかの素粒子よりけた違いに軽いのか、その理由もわかっていません。
\end{phmada}

\begin{phsaki}{相対論と原子核の先へ}
原子核の中の核子の動き（殻模型）、原子核の形や回転、原子核をつくる強い力の理論（量子色力学）は、原子核物理学・素粒子物理学の専門の本で学びます。加速器の設計（加速器工学）では、この章の相対論の運動量とエネルギーを使って、粒の軌道を計算します。重力で時計が進むこと（例題6）や、ブラックホール・重力波は、一般相対性理論で説明されます。
\end{phsaki}

\begin{taisetsu}
\def\tsl#1{\makebox[7.2\zw][l]{\sffamily\bfseries\small\color{\bkTT}#1}}
\tsl{原子核}$A=Z+N$、$R\approx r_0A^{1/3}$（$r_0\approx\qty{1.2}{fm}$）、$\qty{1}{u}=\qty{931.5}{MeV}$\\[3pt]
\tsl{結合エネルギー}$B=(Zm_{\mathrm H}+Nm_{\mathrm n}-M)c^2$\\[3pt]
\tsl{Q 値}$Q=(m_{\text{前}}-m_{\text{後}})c^2$\\[3pt]
\tsl{時間のおくれ}$\Delta t=\gamma\Delta t_0$、$\gamma=\dfrac{1}{\sqrt{1-\beta^2}}$、$L=\dfrac{L_0}{\gamma}$\\[3pt]
\tsl{速度の合成}$u=\dfrac{u'+v}{1+u'v/c^2}$\\[3pt]
\tsl{エネルギー}$p=\gamma mv$、$E=\gamma mc^2$、$E^2=(pc)^2+(mc^2)^2$、$K=(\gamma-1)mc^2$\\[3pt]
\tsl{コンプトン}$\lambda'-\lambda=\dfrac{h}{mc}(1-\cos\theta)$\\[3pt]
\tsl{ドップラー}$f=f_0\sqrt{\dfrac{1-\beta}{1+\beta}}$（遠ざかるとき）
\end{taisetsu}

\begin{kotoba}
\small
\bkword{特殊相対性理論}{光の速さがだれから見ても同じことから、時間と空間の関係を作り直した理論。}
\bkword{静止エネルギー}{止まっている物が質量に応じて持つエネルギー $mc^2$。}
\bkword{同位体}{原子番号が同じで、中性子の数（質量数）が違う原子。}
\bkword{質量欠損}{原子核の質量が、ばらばらの核子の質量の和より軽いぶん。}
\bkword{結合エネルギー}{原子核をばらばらの核子に分けるのに要るエネルギー。質量欠損 $\times c^2$。}
\bkword{Q 値}{核反応の前後の質量の差から出るエネルギー。}
\bkword{核力}{核子どうしを、短い距離で強く引きつける力。}
\bkword{素粒子}{クォークやレプトンなど、物をつくるいちばん小さな粒。}
\end{kotoba}

\bkkigocards{\bkkigocard{$\gamma$}{ガンマ}{ローレンツ因子 $1/\sqrt{1-v^2/c^2}$（単位なし）。原-11 だけの記号。放射線の γ線とは別のもの}{Lorentz factor}{原-11}}
{\bkkigocard{$B$}{ビー}{原子核の結合エネルギー〔MeV〕。原-11 では磁場を出さない。原-1・原-7 の磁場 $\vec B$、原-4 のアインシュタインの係数 $B$ とは別の量}{binding energy}{原-11}}
\bkkigocards{\bkkigocard{u}{ユー}{原子質量単位。炭素12の原子の質量の12分の1。$\qty{1}{u}=\qty{1.66054e-27}{kg}$（931.5 MeV）。立体で書く}{atomic mass unit}{原-11}}
{\bkkigocard{$Q$}{キュー}{核反応で出るエネルギー（Q 値）〔MeV〕。第2巻の熱量 $Q$、第4巻の電荷 $Q$ とは別の量}{Q-value}{原-11}}

%% ======================== 例題 ========================
\Needspace{10\baselineskip}\bksection{reidai}{例題でたしかめよう}
\label{ph:原-11:ex1}
\begin{phreidai}[原子核の密度]
金197 の原子核の半径を求めましょう。また、原子核の密度を求め、1 cm$^3$ の原子核の物質の質量はいくらになるか、計算しましょう。$r_0=\qty{1.2}{fm}$、$\qty{1}{u}=\qty{1.66e-27}{kg}$ とします。
\phbamen
原子核を、核子がぎっしりつまった球とみる場面です（③）。
\phmushi
原子核のふちがぼやけていることと、陽子と中性子の質量の小さな違いは無視します。核子1個の質量を 1 u とします。
\phzu
質量数 $A=197$、半径 $R=r_0A^{1/3}$、質量 $A\times\qty{1}{u}$。
\phriyu
半径は $R\approx r_0A^{1/3}$ で、体積は $A$ に比例します。密度は「質量÷体積」なので、$A$ が消えて、どの原子核でも同じ値になるはずです。
\phtateru
\[ R=r_0A^{1/3},\qquad\rho=\frac{A\times\qty{1}{u}}{\tfrac43\pi R^3}=\frac{\qty{1}{u}}{\tfrac43\pi r_0^{\,3}} \]
\phtoku
$197^{1/3}=5.82$ なので $R=1.2\times5.82=\qty{7.0}{fm}$。
\[ \rho=\frac{1.66\times10^{-27}}{\tfrac43\times3.14\times(1.2\times10^{-15})^3}=\qty{2.3e17}{kg/m^3} \]
1 cm$^3$（$10^{-6}$ m$^3$）では $\qty{2.3e11}{kg}$、つまり約2億3千万トンです。
\phtashikame
単位：kg÷m$^3$＝kg/m$^3$。\\
大きさ：水（$10^3$ kg/m$^3$）の約 $2\times10^{14}$ 倍。原子の中はほとんどからっぽで、質量は原子核に集まっていることがわかります。\\
極端な場合：$A$ を変えても密度は同じ。星が燃えつきてできる中性子星は、星全体がこの密度に近い物質です。
\kotae 半径 約\qty{7.0}{fm}、密度 約 $\qty{2.3e17}{kg/m^3}$、1 cm$^3$ で約 $\qty{2.3e11}{kg}$（約2.3億トン）
\end{phreidai}
\begin{jibun}
鉄56 の原子核の半径を求めましょう。原子の半径を約 \qty{0.1}{nm} とすると、原子核の半径は原子の半径の約何分の1ですか（$r_0=\qty{1.2}{fm}$）。
\bkjans{$56^{1/3}=3.83$ より $R=\qty{4.6}{fm}$。$0.1\times10^{-9}\div(4.6\times10^{-15})$ で、約2万分の1}
\end{jibun}

\label{ph:原-11:ex2}
\begin{phreidai}[結合エネルギーを数える（4けたで書く）]
NIST の原子の質量の表（${}^1$H 1.007825 u、中性子 1.008665 u、${}^{56}$Fe 55.934936 u）から、鉄56 の結合エネルギーと、核子1個あたりの結合エネルギーを求めましょう（表の値と比べるので、4けたで書きます。$\qty{1}{u}=\qty{931.5}{MeV}$）。
\phbamen
鉄56 の原子核を、陽子26個と中性子30個に分けるのに要るエネルギーを求める場面です（④）。
\phmushi
原子の中の電子をつなぎとめるエネルギーは無視します。
\phzu
$Z=26$、$N=56-26=30$。原子の質量 $M=\qty{55.934936}{u}$。
\phriyu
水素原子 26個と中性子 30個の質量の和から、鉄の原子の質量を引くと、質量欠損になります（電子 26個はどちらにも入っていて、消えます）。それに 931.5 をかけます。
\phtateru
\[ \Delta m=Zm_{\mathrm H}+Nm_{\mathrm n}-M,\qquad B=931.5\times\Delta m \]
\phtoku
\[ \Delta m=26\times1.007825+30\times1.008665-55.934936 \]
\[ \phantom{\Delta m}=\qty{0.528464}{u} \]
\[ B=931.5\times0.528464=\qty{492.3}{MeV} \]
\[ \frac{B}{A}=\frac{492.3}{56}=\qty{8.790}{MeV} \]
\phtashikame
単位：u×（MeV/u）＝MeV。\\
大きさ：図\ref{ph:原-11:fig-curve}の山の高さ（約 8.8 MeV）と合います。ヘリウム4 の 7.07 MeV より大きく、鉄はとても固い原子核です。\\
極端な場合：核子が1個だけの水素では、質量欠損は0です。
\kotae $B=\qty{492.3}{MeV}$、1核子あたり \qty{8.790}{MeV}
\end{phreidai}
\begin{jibun}
炭素12（原子の質量は 12 u ちょうど）の結合エネルギーと、核子1個あたりの結合エネルギーを求めましょう（${}^1$H 1.007825 u、中性子 1.008665 u、1 u＝931.5 MeV）。
\bkjans{$\Delta m=6\times1.007825+6\times1.008665-12=\qty{0.098940}{u}$、$B=\qty{92.16}{MeV}$、1核子あたり \qty{7.680}{MeV}}
\end{jibun}

\label{ph:原-11:ex3}
\begin{phreidai}[核融合と核分裂の Q 値]
（1）重水素と三重水素の核融合 ${}^2_1\mathrm H+{}^3_1\mathrm H\to{}^4_2\mathrm{He}+{}^1_0\mathrm n$ の Q 値を求めましょう。（2）ウラン235 に中性子が当たって分かれる反応 ${}^{235}_{92}\mathrm U+{}^1_0\mathrm n\to{}^{141}_{56}\mathrm{Ba}+{}^{92}_{36}\mathrm{Kr}+3\,{}^1_0\mathrm n$ の Q 値を求めましょう。原子の質量は ${}^2$H 2.014102、${}^3$H 3.016049、${}^4$He 4.002603、n 1.008665、${}^{235}$U 235.043930、${}^{141}$Ba 140.914403、${}^{92}$Kr 91.926173（どれも u）、1 u＝931.5 MeV とします。
\phbamen
核融合と核分裂で、1回の反応で出るエネルギーを比べる場面です（⑥）。
\phmushi
ぶつける前の粒の運動エネルギーは、出るエネルギーより十分小さいとして無視します。
\phzu
左下（電荷）：（1）$1+1=2+0$、（2）$92+0=56+36+0$。左上（核子の数）：（1）$2+3=4+1$、（2）$235+1=141+92+3$。どちらも合っています。電子の数も左右で同じです。
\phriyu
Q 値は、反応の前と後の質量の差に 931.5 をかけたものです。
\phtateru
\[ Q=931.5\times(m_{\text{前}}-m_{\text{後}}) \]
\phtoku
（1）
\[ m_{\text{前}}-m_{\text{後}}=(2.014102+3.016049)-(4.002603+1.008665)=\qty{0.018883}{u} \]
\[ Q=931.5\times0.018883=\qty{17.6}{MeV} \]
（2）
\[ (235.043930+1.008665)-(140.914403+91.926173+3\times1.008665) \]
\[ =\qty{0.186024}{u},\qquad Q=931.5\times0.186024=\qty{173}{MeV} \]
\phtashikame
単位：MeV。\\
大きさ：核分裂のほうが1回のエネルギーは約10倍大きいが、反応する質量でわると、核融合は質量の 0.38\,\%、核分裂は 0.079\,\% が熱になり、核融合のほうが約5倍「質量あたりのエネルギー」が大きい。\\
極端な場合：（2）では、このあとバリウムやクリプトンがβ崩壊を繰り返し、さらにエネルギーが出るので、合わせて約 \qty{200}{MeV} になります（原-15）。
\kotae （1）\qty{17.6}{MeV}　（2）\qty{173}{MeV}
\end{phreidai}
\begin{jibun}
重水素どうしの核融合 ${}^2_1\mathrm H+{}^2_1\mathrm H\to{}^3_2\mathrm{He}+{}^1_0\mathrm n$ の Q 値を求めましょう（${}^2$H 2.014102 u、${}^3$He 3.016029 u、n 1.008665 u、1 u＝931.5 MeV）。
\bkjans{$2\times2.014102-3.016029-1.008665=\qty{0.003510}{u}$、$Q=\qty{3.27}{MeV}$}
\end{jibun}

\label{ph:原-11:ex4}
\begin{phreidai}[太陽が1秒に失う質量（はてな？）]
太陽の中心では、水素の原子核4個から、ヘリウム4 の原子核が1個できます。原子の質量で ${}^1$H 1.007825 u、${}^4$He 4.002603 u として、（1）1回の反応で出るエネルギー、（2）太陽が1秒に失う質量、（3）1秒にヘリウムに変わる水素の質量を求めましょう。太陽の明るさは $\qty{3.83e26}{W}$、1 u＝931.5 MeV、$c=\qty{3.00e8}{m/s}$ とします。
\phbamen
太陽の光のエネルギーが、どれだけの質量から生まれるかを調べる場面です（②④⑪）。
\phmushi
ニュートリノが持ち去るエネルギー（数\,\%）と、太陽の中の反応の細かい道すじは無視します。陽電子は電子と消え合ってγ線になり、熱になるとします。
\phzu
反応：$4\,{}^1\mathrm H\to{}^4\mathrm{He}$（原子どうし）。明るさ $L=\qty{3.83e26}{W}$。
\phriyu
質量の差に 931.5 をかけると1回のエネルギーです。太陽が1秒に失う質量は、出したエネルギーを $c^2$ でわれば求まります（$E=mc^2$）。
\phtateru
\[ Q=931.5\times(4m_{\mathrm H}-m_{\mathrm{He}}),\qquad\Delta m=\frac{L\times\qty{1}{s}}{c^2} \]
\phtoku
（1）$4\times1.007825-4.002603=\qty{0.028697}{u}$、$Q=\qty{26.7}{MeV}$。水素4個の質量の $0.028697\div4.031300=0.712\,\%$ です。\\
（2）1秒に減る質量は
\[ \Delta m=\frac{3.83\times10^{26}}{(3.00\times10^8)^2}=\qty{4.3e9}{kg} \]
（3）失う質量は、変わる水素の 0.712\,\% なので、$4.3\times10^9\div0.00712=\qty{6.0e11}{kg}$。
\phtashikame
単位：J÷(m/s)$^2$＝kg。\\
大きさ：1秒に約43億 kg ずつ軽くなっていますが、太陽の質量（$\qty{2.0e30}{kg}$）に比べると、46億年で 0.03\,\% ほどです。\\
極端な場合：反応が化学の反応（質量の $10^{-10}$ ほど）なら、同じ明るさに必要な燃料は何千万倍にもなります。
\kotae （1）\qty{26.7}{MeV}（0.71\,\%）　（2）約 $\qty{4.3e9}{kg}$　（3）約 $\qty{6.0e11}{kg}$
\end{phreidai}
\begin{jibun}
水素 \qty{1}{kg} が全部ヘリウムに変わると、何 J のエネルギーが出ますか。質量の 0.712\,\% がエネルギーになるとし、$c=\qty{3.00e8}{m/s}$ とします。
\bkjans{$0.00712\times1\times(3.00\times10^8)^2=\qty{6.4e14}{J}$}
\end{jibun}

\label{ph:原-11:ex5}
\begin{phreidai}[地上に届くミューオン]
宇宙から来た粒が、高さ十数 km の大気にぶつかって、ミューオンが生まれます。ミューオンは、止まっていれば平均 \qty{2.2}{\micro s} でこわれます。速さ $0.995c$ のミューオンについて、（1）$\gamma$、（2）地上から見た平均の寿命、（3）こわれるまでに進む平均の距離を、時間のおくれを考えないときと、考えたときで求めましょう（$c=\qty{3.00e8}{m/s}$）。
\phbamen
光の速さに近い粒の時計がおくれる場面です（⑦）。
\phmushi
ミューオンが空気の中で遅くなること、地球の重力は無視します。速さは一定とします。
\phzu
$\beta=0.995$、止まっているときの平均の寿命 $\Delta t_0=\qty{2.2}{\micro s}$。
\phriyu
地上の人から見ると、動いているミューオンの時計は $\gamma$ 倍ゆっくり進むので、寿命は $\gamma\Delta t_0$ になります。
\phtateru
\[ \gamma=\frac{1}{\sqrt{1-\beta^2}},\qquad\Delta t=\gamma\Delta t_0,\qquad d=v\,\Delta t \]
\phtoku
（1）$1-0.995^2=0.009975$、$\sqrt{0.009975}=0.0999$ なので $\gamma=10.0$。\\
（2）$\Delta t=10.0\times2.2=\qty{22}{\micro s}$。\\
（3）おくれを考えないと $0.995\times3.00\times10^8\times2.2\times10^{-6}=\qty{660}{m}$。考えると、その10倍の約 \qty{6.6}{km}。
\phtashikame
単位：m/s×s＝m。\\
大きさ：おくれを考えないと、ほとんどのミューオンは地上に届かないはずですが、実際には地上でたくさん見つかります。時間のおくれの証拠です。\\
極端な場合：ミューオンと動く人から見ると、寿命は \qty{2.2}{\micro s} のままで、大気の厚さが10分の1に縮んで見えます（長さの縮み）。
\kotae （1）$\gamma=10$　（2）約 \qty{22}{\micro s}　（3）約 \qty{0.66}{km} と、約 \qty{6.6}{km}
\end{phreidai}
\begin{jibun}
速さ $0.98c$ のミューオンの $\gamma$、地上から見た平均の寿命、こわれるまでに進む平均の距離を求めましょう（止まっているときの寿命 \qty{2.2}{\micro s}、$c=\qty{3.00e8}{m/s}$）。
\bkjans{$\gamma=1/\sqrt{1-0.98^2}=5.0$、寿命 約\qty{11}{\micro s}、距離 $0.98\times3.00\times10^8\times11\times10^{-6}$ で約\qty{3.2}{km}}
\end{jibun}

\label{ph:原-11:ex6}
\begin{phreidai}[GPS の衛星の時計]
GPS の衛星は、地球の中心から $r=\qty{2.656e7}{m}$ の円軌道を回っています。（1）衛星の速さ $v$ を求め、（2）速さによる時計のおくれが1日に何 \si{\micro s} か求めましょう。（3）重力が弱いことによる時計の進みは、一般相対性理論から1日に \qty{45.7}{\micro s} です。合わせると、衛星の時計は1日に何 \si{\micro s} 進みますか。その時間に電波が進む距離も求めましょう。地球の $GM=\qty{3.986e14}{m^3/s^2}$、$c=\qty{3.00e8}{m/s}$ とします。
\phbamen
相対論の時間のおくれが、身近な道具に効く場面です（⑦。物理シリーズの『波動』（波-14））。
\phmushi
地球の自転と、軌道のわずかなだ円は無視します。
\phzu
衛星の速さ $v$、地上の1日 $\qty{86400}{s}$。
\phriyu
円軌道の速さは、万有引力が向心力になる条件 $GMm/r^2=mv^2/r$（力-6）から求まります。$\beta$ はとても小さいので、$\gamma\approx1+\tfrac12\beta^2$ で、1秒あたりのおくれは $\tfrac12\beta^2$ 秒です。
\phtateru
\[ v=\sqrt{\frac{GM}{r}},\qquad\text{1日のおくれ}=\frac12\frac{v^2}{c^2}\times\qty{86400}{s} \]
\phtoku
（1）$v=\sqrt{3.986\times10^{14}\div(2.656\times10^7)}=\qty{3.87e3}{m/s}$\\
（2）1日のおくれは
\[ \frac12\times\frac{(3.87\times10^3)^2}{(3.00\times10^8)^2}\times86400=\qty{7.2e-6}{s} \]
つまり1日に約 \qty{7.2}{\micro s} おくれる。\\
（3）$45.7-7.2=\qty{38.5}{\micro s}$ 進む。電波が進む距離は
\[ 3.00\times10^8\times38.5\times10^{-6}=\qty{1.2e4}{m} \]
\phtashikame
単位：(m/s)$^2$÷(m/s)$^2$×s＝s。\\
大きさ：力-1 と波-14 で紹介した「1日に約 \qty{38}{\micro s}」と合います。\\
極端な場合：補正しなければ、1日で約 \qty{12}{km} 分の時刻のずれがたまり、位置が大きくずれます。
\kotae （1）約 $\qty{3.9e3}{m/s}$　（2）約 \qty{7.2}{\micro s}　（3）約 \qty{38.5}{\micro s} 進む。約 \qty{12}{km}
\end{phreidai}
\begin{jibun}
国際宇宙ステーションは、秒速約 \qty{7.7}{km} で地球を回っています。速さによる時計のおくれは、1日に約何 \si{\micro s} ですか（$c=\qty{3.00e8}{m/s}$。重力の効果は考えない）。
\bkjans{$\tfrac12\times(7.7\times10^3)^2$ $\div(3.00\times10^8)^2$ $\times86400$ で、約\qty{28}{\micro s}}
\end{jibun}

\label{ph:原-11:ex7}
\begin{phreidai}[速度の合成（3けたで書く）]
光の速さの0.6倍で飛ぶロケットから、前へ、ロケットから見て光の速さの0.6倍で小さな探査機を打ち出しました。地面から見た探査機の速さを求めましょう（力-1 の約束。3けたで書きます）。
\phbamen
光の速さに近い速度を合成する場面です（⑦）。
\phmushi
ロケットと探査機は、一直線の上を一定の速さで動くとします。
\phzu
$v=0.6c$（地面から見たロケット）、$u'=0.6c$（ロケットから見た探査機）。
\phriyu
光の速さに近いので、$u'+v$ ではなく、ローレンツ変換から導いた合成の式を使います。
\phtateru
\[ u=\frac{u'+v}{1+u'v/c^2} \]
\phtoku
\[ u=\frac{0.6c+0.6c}{1+0.6\times0.6}=\frac{1.2c}{1.36}=0.882c \]
\phtashikame
単位：両辺とも速さ。\\
大きさ：$c$ をこえません。ふつうの足し算の $1.2c$ よりずっと小さい。\\
極端な場合：力-1 の $0.5c+0.5c$ では $1/1.25=0.8c$。時速100 km どうしでは、ずれは $10^{-13}$ m/s ほどで、まったく気づきません。
\kotae $0.882c$
\end{phreidai}
\begin{jibun}
（1）$0.9c$ で飛ぶロケットから、前へ $0.9c$（ロケットから見て）で探査機を打ち出すと、地面から見た速さはいくらですか。（2）同じロケットから前へ光を出すと、地面から見た光の速さはいくらですか。
\bkjans{（1）$1.8c\div1.81=0.994c$　（2）$(c+0.9c)\div(1+0.9)=c$}
\end{jibun}

\label{ph:原-11:ex8}
\begin{phreidai}[電圧で加速した電子の速さ]
止まっていた電子を、（1）\qty{1}{kV}、（2）\qty{100}{kV} の電圧で加速します。それぞれの速さは光の速さの何倍ですか。ニュートンの式で求めた値と比べましょう（$m_{\mathrm e}c^2=\qty{0.511}{MeV}$）。
\phbamen
電子顕微鏡（原-3）や X線の管で、電子を高い電圧で加速する場面です（⑧。電-2 の約束）。
\phmushi
加速する前の速さと、電子が出す光（電-15）は無視します。
\phzu
運動エネルギー（1）$K=\qty{1}{keV}=\qty{0.001}{MeV}$、（2）$K=\qty{0.1}{MeV}$。
\phriyu
$K=(\gamma-1)mc^2$ から $\gamma$ が求まり、$\gamma$ から $\beta=\sqrt{1-1/\gamma^2}$ です。ニュートンの式では $\tfrac12mv^2=K$ から $\beta=\sqrt{2K/(mc^2)}$ です。
\phtateru
\[ \gamma=1+\frac{K}{mc^2},\qquad\beta=\sqrt{1-\frac{1}{\gamma^2}} \]
\phtoku
（1）$\gamma=1+0.001/0.511=1.00196$、$\beta=\sqrt{1-1/1.00196^2}=0.0625$。ニュートンの式でも 0.0626 で、ほとんど同じ。\\
（2）$\gamma=1+0.1/0.511=1.196$、$\beta=\sqrt{1-1/1.196^2}=0.548$。ニュートンの式では 0.626 で、14\,\% も大きくなってしまう。
\phtashikame
単位：$\beta$ は単位なし。\\
大きさ：（1）は電-2 例題5 の「光の速さの約6\,\%」と合います。\\
極端な場合：\qty{1}{MV} では、ニュートンの式は $1.98c$ と光をこえてしまいますが、相対論では $0.941c$ です。
\kotae （1）$0.0625c$（ニュートンの式とほぼ同じ）　（2）$0.548c$（ニュートンの式では $0.626c$）
\end{phreidai}
\begin{jibun}
止まっていた電子を \qty{1}{MV} で加速すると、速さは光の速さの何倍ですか。ニュートンの式ではどうなりますか（$m_{\mathrm e}c^2=\qty{0.511}{MeV}$）。
\bkjans{$\gamma=1+1/0.511=2.96$、$\beta=0.941$。ニュートンの式では $\sqrt{2\times1/0.511}=1.98$ で、光より速くなってしまう}
\end{jibun}

\label{ph:原-11:ex9}
\begin{phreidai}[速い電子の運動量]
運動エネルギーが \qty{1.00}{MeV} の電子の、全部のエネルギー $E$ と運動量 $p$ を求めましょう（$m_{\mathrm e}c^2=\qty{0.511}{MeV}$、$\qty{1}{MeV}=\qty{1.60e-13}{J}$、$c=\qty{3.00e8}{m/s}$）。
\phbamen
放射線のβ線（原-12）のような、光の速さに近い電子の運動量を求める場面です（⑧）。
\phmushi
とくになし（1個の自由な電子とする）。
\phzu
$K=\qty{1.00}{MeV}$、$mc^2=\qty{0.511}{MeV}$。
\phriyu
$E=K+mc^2$ で、$E^2=(pc)^2+(mc^2)^2$ から $pc$ を求めます。MeV で計算してから、最後に kg·m/s に直します。
\phtateru
\[ E=K+mc^2,\qquad pc=\sqrt{E^2-(mc^2)^2} \]
\phtoku
$E=1.00+0.511=\qty{1.511}{MeV}$、$pc=\sqrt{1.511^2-0.511^2}=\qty{1.42}{MeV}$。
\[ p=\frac{1.42\times1.60\times10^{-13}}{3.00\times10^8}=\qty{7.6e-22}{kg.m/s} \]
\phtashikame
単位：MeV÷(m/s) は J·s/m＝kg·m/s。\\
大きさ：ニュートンの式 $p=\sqrt{2mK}$ では $pc=\qty{1.01}{MeV}$ で、約3割小さくなってしまいます。\\
極端な場合：$K$ がとても大きいと、$pc\approx E$ で、光子と同じ関係になります（図\ref{ph:原-11:fig-Ep}）。
\kotae $E=\qty{1.51}{MeV}$、$pc=\qty{1.42}{MeV}$（$p\approx\qty{7.6e-22}{kg.m/s}$）
\end{phreidai}
\begin{jibun}
運動エネルギーが \qty{1000}{MeV} の陽子の $pc$ を求めましょう。速さは光の速さの何倍ですか（$m_{\mathrm p}c^2=\qty{938}{MeV}$）。
\bkjans{$E=\qty{1938}{MeV}$、$pc=\sqrt{1938^2-938^2}\approx\qty{1696}{MeV}$、$\beta=pc/E=0.875$}
\end{jibun}

\label{ph:原-11:ex10}
\begin{phreidai}[コンプトン散乱でエネルギーを失う γ線]
セシウム137 から出る \qty{0.662}{MeV} の γ線が、物の中の電子にぶつかって、90° の向きに散らばりました。（1）波長ののび、（2）散らばった γ線のエネルギー、（3）電子が受けとった運動エネルギーを求めましょう（$h/(m_{\mathrm e}c)=\qty{2.43}{pm}$、$m_{\mathrm e}c^2=\qty{0.511}{MeV}$）。
\phbamen
放射線が物にあたって、エネルギーを渡す場面です（⑧。原-13 で詳しく学ぶ）。
\phmushi
電子を原子につなぎとめているエネルギー（数 eV〜数十 keV）は、\qty{0.662}{MeV} より十分小さいので、止まった自由な電子とみなします。
\phzu
図\ref{ph:原-11:fig-compton}。$E=\qty{0.662}{MeV}$、$\theta=90^\circ$。
\phriyu
⑧で導いたコンプトンの式を、エネルギーの形 $\dfrac1{E'}-\dfrac1E=\dfrac{1-\cos\theta}{mc^2}$ で使います。
\phtateru
\[ \Delta\lambda=\frac{h}{mc}(1-\cos\theta),\qquad E'=\frac{E}{1+\dfrac{E}{mc^2}(1-\cos\theta)} \]
\phtoku
（1）$\cos90^\circ=0$ なので $\Delta\lambda=\qty{2.43}{pm}$。\\
（2）$E'=\dfrac{0.662}{1+0.662/0.511}=\dfrac{0.662}{2.296}=\qty{0.288}{MeV}$\\
（3）$0.662-0.288=\qty{0.374}{MeV}$
\phtashikame
単位：MeV。\\
大きさ：もとの γ線の波長は $1240/(0.662\times10^6)=\qty{1.87}{pm}$ で、2.43 pm のびて 4.30 pm になります。$hc/E'$ で求めた $1240/(0.288\times10^6)=\qty{4.31}{pm}$ と合います。（2）（3）は、原-2 例題10 で式を読んで求めた値と同じです。\\
極端な場合：$\theta=0$ なら、波長は変わらず、電子はエネルギーを受けとりません。
\kotae （1）\qty{2.43}{pm}　（2）\qty{0.288}{MeV}　（3）\qty{0.374}{MeV}
\end{phreidai}
\begin{jibun}
例題10 の γ線が、真後ろ（180°）にはね返されたとき、波長ののびと、はね返された γ線のエネルギーを求めましょう。
\bkjans{$\Delta\lambda=2\times2.43=\qty{4.85}{pm}$、$E'=0.662\div(1+2\times0.662/0.511)=\qty{0.184}{MeV}$}
\end{jibun}

\label{ph:原-11:ex11}
\begin{phreidai}[遠ざかる星の光（3けたで書く）]
光の速さの0.1倍で遠ざかる天体から来る光の振動数は、もとの振動数の何倍ですか。音のドップラー効果の式（見る人が止まり、音源が遠ざかる）と、近似の式 $f\approx f_0(1-\beta)$ の値と比べましょう（3けたで書きます）。
\phbamen
光源が見る人から遠ざかる場面です（⑨。波-8 の約束）。
\phmushi
天体は見る人と結ぶ線の上をまっすぐ遠ざかるとし、重力の効果や宇宙の膨張は無視します。
\phzu
$\beta=0.1$。
\phriyu
光には媒質がないので、⑨で導いた相対論の式を使います。
\phtateru
\[ \frac{f}{f_0}=\sqrt{\frac{1-\beta}{1+\beta}} \]
\phtoku
\[ \frac{f}{f_0}=\sqrt{\frac{0.9}{1.1}}=\sqrt{0.8182}=0.905 \]
音の式は $1/(1+\beta)=0.909$、近似は $1-\beta=0.900$ です。
\phtashikame
単位：比なので単位なし。\\
大きさ：3つの値は1\,\% 以内で合います。$\beta=0.1$ くらいまでなら、波-8 の近似で十分です。\\
極端な場合：$\beta\to1$ で $f\to0$（光が届かなくなる）。水素の 656.3 nm の線は、725.6 nm に見えます。
\kotae 0.905 倍（音の式 0.909、近似 0.900）
\end{phreidai}
\begin{jibun}
水素が出す \qty{656.3}{nm} の光が、ある星からは \qty{663.0}{nm} に見えました。星は近づいていますか、遠ざかっていますか。速さを求めましょう（$\lambda=c/f$ なので、遠ざかるとき $\lambda=\lambda_0\sqrt{(1+\beta)/(1-\beta)}$）。
\bkjans{波長がのびているので遠ざかる。$(663.0/656.3)^2=(1+\beta)/(1-\beta)$ より $\beta=0.0102$、速さ 約 $\qty{3.0e6}{m/s}$}
\end{jibun}

\label{ph:原-11:ex12}
\begin{phreidai}[原子核の中の陽子のクーロン力と、核力の届く距離]
（1）陽子どうしが \qty{1.0}{fm} はなれているときのクーロン力と、クーロン力による位置エネルギー（MeV）を求めましょう（$k=\qty{8.99e9}{N.m^2/C^2}$、$e=\qty{1.60e-19}{C}$）。（2）π中間子の静止エネルギーを \qty{140}{MeV} として、湯川の考えで、核力の届く距離 $\hbar c/(mc^2)$ を求めましょう（$\hbar c=\qty{197.3}{MeV.fm}$）。
\phbamen
原子核の中で、陽子どうしが押し合う力と、それに勝つ核力の大きさを見積もる場面です（⑩。電-1 の約束）。
\phmushi
原子核の中のほかの陽子・中性子の影響は無視し、2個だけを考えます。
\phzu
$r=\qty{1.0e-15}{m}$、電荷はどちらも $e$。
\phriyu
クーロン力は $F=ke^2/r^2$、位置エネルギーは $U=ke^2/r$（電-2）です。力を運ぶ粒の質量から力の届く距離を見積もる湯川の考えを使います。
\phtateru
\[ F=\frac{ke^2}{r^2},\qquad U=\frac{ke^2}{r},\qquad d\approx\frac{\hbar c}{mc^2} \]
\phtoku
（1）$F=\dfrac{8.99\times10^9\times(1.60\times10^{-19})^2}{(1.0\times10^{-15})^2}=\qty{230}{N}$\\
$U=\dfrac{8.99\times10^9\times(1.60\times10^{-19})^2}{1.0\times10^{-15}}=\qty{2.30e-13}{J}=\qty{1.44}{MeV}$\\
（2）$d=197.3\div140=\qty{1.4}{fm}$
\phtashikame
単位：N・J、MeV·fm÷MeV＝fm。\\
大きさ：230 N は、約 \qty{23}{kg} の物の重さです。たった2つの陽子にこの力がはたらいていて、それに勝つ核力は、とても強い力です。位置エネルギー 1.4 MeV は、核子1個あたりの結合エネルギー（約 8 MeV）より小さい。\\
極端な場合：$r$ を2倍にすると、力は4分の1、位置エネルギーは2分の1。
\kotae （1）約 \qty{230}{N}、約 \qty{1.4}{MeV}　（2）約 \qty{1.4}{fm}
\end{phreidai}
\begin{jibun}
陽子どうしが \qty{2.0}{fm} はなれているときのクーロン力と位置エネルギー（MeV）を求めましょう。
\bkjans{$F=230\div4\approx\qty{58}{N}$、$U=1.44\div2=\qty{0.72}{MeV}$}
\end{jibun}

\label{ph:原-11:ex13}
\begin{phreidai}[半経験的質量公式で、鉄とウランを計算する]
半経験的質量公式（係数は⑩の値）で、鉄56（$Z=26$）とウラン238（$Z=92$）の結合エネルギーを計算し、NIST の質量から求めた値（鉄 \qty{492.3}{MeV}、ウラン \qty{1801.7}{MeV}）と比べましょう。どの項が大きく効いていますか。
\phbamen
原子核を電気を帯びたしずくとみて、結合エネルギーを見積もる場面です（⑩）。
\phmushi
殻の構造（魔法数）と、原子核の形のゆがみは無視します。
\phzu
鉄：$A=56$、$Z=26$、$N=30$（偶偶）。ウラン：$A=238$、$Z=92$、$N=146$（偶偶）。
\phriyu
5つの項を、1つずつ計算して足します。どちらも陽子・中性子が偶数なので、対の項は $+12.0/\sqrt A$ です。
\phtateru
\[ B=15.8A-18.3A^{2/3}-0.714\frac{Z(Z-1)}{A^{1/3}}-23.2\frac{(A-2Z)^2}{A}+\frac{12.0}{\sqrt A} \]
\phtoku
\begin{bkcenter}
\begin{bktable}{colspec={Q[l,m,wd=19mm] Q[r,m,wd=14mm] Q[r,m,wd=14mm] Q[r,m,wd=15mm] Q[r,m,wd=13mm] Q[r,m,wd=8mm] Q[r,m,wd=14mm]}, colsep=3pt, rowsep=0.5pt, row{1}={bg=bkTheme!10}, cells={font=\small}}
 & 体積 & 表面 & クーロン & 対称 & 対 & 合計 \\
鉄56 & 884.8 & $-267.9$ & $-121.3$ & $-6.6$ & 1.6 & 490.6 \\
ウラン238 & 3760.4 & $-702.8$ & $-964.6$ & $-284.2$ & 0.8 & 1809.6 \\
\end{bktable}
\end{bkcenter}
\phtashikame
単位：MeV。\\
大きさ：鉄は 490.6（本当は 492.3）、ウランは 1809.6（本当は 1801.7）で、どちらも 0.5\,\% ほどのずれです。鉄では表面の項、ウランではクーロンの項が大きく効いています。\\
極端な場合：陽子がもっと多い原子核では、クーロンの項がさらに大きくなり、ばらばらになろうとする（分裂しやすい）ことがわかります。
\kotae 鉄 約 \qty{491}{MeV}、ウラン 約 \qty{1810}{MeV}（どちらも本当の値と 0.5\,\% ほどのずれ）。鉄は表面の項、ウランはクーロンの項が大きい
\end{phreidai}
\begin{jibun}
半経験的質量公式で、酸素16（$Z=8$）の結合エネルギーを計算し、本当の値 \qty{127.6}{MeV} と比べましょう。
\bkjans{体積 252.8、表面 $-116.2$、クーロン $-15.9$、対称 0、対 $+3.0$ で、約 \qty{123.7}{MeV}。本当の値より約 \qty{4}{MeV} 小さい（酸素16 は陽子・中性子とも8個の魔法数で、式より強く結びついている）}
\end{jibun}

%% ======================== まちがえやすいところ ========================
\Needspace{11\baselineskip}\bksection{tsumazuki}{まちがえやすいところをチェック}
\begin{tsumazuki}{質量がエネルギーに変わるのは、原子核の反応だけ？}
\bkng{化学の反応では、質量は変わらない。$E=mc^2$ は原子核の反応だけの式。}
\bkok{化学の反応でも、出たエネルギーの分だけ質量は減る。ただ、割合が $10^{-10}$ ほどと小さすぎてはかれない。}
$E=mc^2$ は、どんなエネルギーの出入りにも成り立ちます（②）。
\end{tsumazuki}

\begin{tsumazuki}{結合エネルギーが大きい原子核は、重い？}
\bkng{結合エネルギーが大きいほど、エネルギーをたくさん持っているので、原子核は重い。}
\bkok{結合エネルギーは、くっつくときに\textbf{出ていった}エネルギー。大きいほど、ばらばらのときより軽い。}
結合エネルギーは、原子核を分けるのに「外から加える」エネルギーです（④）。
\end{tsumazuki}

\begin{tsumazuki}{核反応の式は、左上の数がそろえばよい？}
\bkng{${}^{14}\mathrm N+{}^4\mathrm{He}\to{}^{17}\mathrm O+{}^1\mathrm n$ は、左上が $14+4=17+1$ で合っている。}
\bkok{左下（電荷）も、$7+2=8+1$ でなければならない。中性子（左下 0）では合わず、陽子 ${}^1_1\mathrm p$ が正しい。}
核子の数と電荷の、2つの保存を確かめます（⑥）。
\end{tsumazuki}

\begin{tsumazuki}{速度は、いつも足し算？}
\bkng{$0.6c$ のロケットから $0.6c$ で打ち出すと、地面からは $1.2c$。}
\bkok{$u=(u'+v)/(1+u'v/c^2)=0.882c$。どんな速度を合成しても $c$ をこえない。}
足し算でよいのは、$c$ よりずっと遅いときだけです（⑦）。
\end{tsumazuki}

\begin{tsumazuki}{速く動くと、質量が増える？}
\bkng{光の速さに近づくと、物の質量が $\gamma m$ に増えるので、加速しにくくなる。}
\bkok{質量 $m$ はどの慣性系でも同じ。$\gamma$ 倍になるのは運動量 $p=\gamma mv$ とエネルギー $E=\gamma mc^2$。}
古い本には「相対論的質量 $\gamma m$」と書いたものもありますが、いまの専門の本では、質量は $E^2=(pc)^2+(mc^2)^2$ の $m$ だけを指します（⑧）。
\end{tsumazuki}

\begin{tsumazuki}{動いている時計がおくれるのは、時計がこわれるから？}
\bkng{速く動くと、時計のしくみが乱れて、おくれる。}
\bkok{どんな時計も同じ割合でおくれる。時間そのものの進み方が、見る人によって違う。}
光の時計でも、原子の時計でも、ミューオンの寿命でも、同じ $\gamma$ でおくれます（⑦）。
\end{tsumazuki}

\begin{column}{1905年のアインシュタイン}
アインシュタインは、スイスの特許局で働いていた26歳の1905年に、物理を大きく変える論文を次々に発表しました。光は粒（光子）としてふるまうという光電効果の論文（原-2。1921年のノーベル物理学賞の理由になった）、水にうかぶ花粉の小さな粒が不規則に動くわけを、水の分子の衝突から説明した論文、特殊相対性理論の論文、そして「物体の慣性は、そのエネルギーによるか」という短い論文です。最後の論文で、物がエネルギーを出すと質量が減ること、つまり $E=mc^2$ の考えが初めて示されました。原子核のエネルギーが確かめられるのは、それから約30年後のことです（⑥）。
\end{column}

%% ======================== 工学ではこう使う ========================
\Needspace{18\baselineskip}\bksection{phgear}{工学ではこう使う}
\begin{phkougaku}[質量とエネルギーを数える技術]
\textbf{原子力と核融合}：原子炉は、ウランの核分裂で出る約 \qty{200}{MeV} を熱にして、発電に使います（原-15）。核融合炉の研究では、D と T の反応の \qty{17.6}{MeV} のうち、中性子が持ち出す \qty{14.1}{MeV} を、炉の壁で熱に変える設計をしています。

\textbf{加速器と医療}：病院の放射線治療の装置は、電子を数 MeV〜20 MeV ほどに加速して、X線を作ります。電子の速さは光の速さの 99\,\% をこえ、装置は相対論の運動量で設計されています。PET 検査では、陽電子と電子が消えてできる \qty{0.511}{MeV} の γ線2本をとらえます。

\textbf{時間とナビゲーション}：GPS の衛星の時計は、相対論の補正を前もって入れて打ち上げられます（例題6）。

\begin{phketa}[原子核と相対論の数の目安]
\phketaline{電子・陽子の静止エネルギー}{\qty{0.511}{MeV}・\qty{938}{MeV}}
\phketaline{核子1個あたりの結合エネルギー}{約 \qty{8}{MeV}（鉄で 8.8 MeV）}
\phketaline{核分裂・D-T 核融合の1回}{約 \qty{200}{MeV}・\qty{17.6}{MeV}}
\phketaline{GPS の時計のずれ}{1日に約 \qty{38}{\micro s}}
\phboxhead{答えが桁で合っているか}
原子核の反応のエネルギーは、ほとんどが 1〜200 MeV のあいだです。核子1個あたりにすると、1 MeV 前後から 10 MeV ほど。何 GeV や何 eV になったら、u と MeV の換算（931.5）や、$c^2$ のかけ忘れを疑いましょう。
\end{phketa}
\end{phkougaku}
\phstation{原子核工学,加速器工学,医学物理,宇宙工学}{質量欠損、結合エネルギー、Q 値、静止エネルギー、ローレンツ因子、相対論的運動量、コンプトン散乱}

%% ======================== 章末問題 ========================
\bkendproblems
\bkproblemgroup{基本}{3}
\bkproblem ネオン20（陽子・中性子10個ずつ、19.992440 u）とケイ素28（14個ずつ、27.976927 u）の結合エネルギーと、核子1個あたりの値を求めましょう。どちらの核子が強く結びついていますか（${}^1$H 1.007825 u、中性子 1.008665 u、1 u＝931.5 MeV）。
\bkpans{ネオン20：$\Delta m=10\times1.007825$ $+10\times1.008665$ $-19.992440$ $=\qty{0.172460}{u}$、$B\approx\qty{160.6}{MeV}$、1核子あたり \qty{8.03}{MeV}。ケイ素28：$\Delta m=\qty{0.253933}{u}$、$B\approx\qty{236.5}{MeV}$、1核子あたり \qty{8.45}{MeV}。ケイ素28 のほうが強い（鉄に近いほど大きい）}
\bkproblem 核融合炉で三重水素をつくる反応 ${}^6_3\mathrm{Li}+{}^1_0\mathrm n\to{}^3_1\mathrm H+{}^4_2\mathrm{He}$ の Q 値を求めましょう（${}^6$Li 6.015123 u、中性子 1.008665 u、${}^3$H 3.016049 u、${}^4$He 4.002603 u、1 u＝931.5 MeV）。
\bkpans{$6.015123+1.008665$ $-3.016049-4.002603$ $=\qty{0.005136}{u}$、$Q\approx\qty{4.78}{MeV}$}
\bkproblem 次の核反応式の □ に入る粒（または原子核）を書きましょう。（1）${}^{235}_{92}\mathrm U\to\square+{}^4_2\mathrm{He}$（原子番号 90 の元素はトリウム Th）　（2）${}^{27}_{13}\mathrm{Al}+{}^4_2\mathrm{He}\to\square+{}^1_0\mathrm n$　（3）${}^{14}_7\mathrm N+{}^1_0\mathrm n\to{}^{14}_6\mathrm C+\square$
\bkpans{（1）トリウム231 ${}^{231}_{90}\mathrm{Th}$　（2）リン30 ${}^{30}_{15}\mathrm P$　（3）陽子 ${}^1_1\mathrm p$（大気の上のほうで、宇宙線の中性子が炭素14 をつくる反応）}
\bkproblem 鉛208 とヘリウム4 の原子核の半径を求めましょう。半径・体積は何倍違いますか（$r_0=\qty{1.2}{fm}$）。
\bkpans{鉛 約\qty{7.1}{fm}、ヘリウム 約\qty{1.9}{fm}。半径は $(208/4)^{1/3}=52^{1/3}\approx3.7$ 倍、体積は質量数に比例して52倍}
\bkproblem \qty{0.1}{g} の質量が全部エネルギーに変わると何 J ですか。\qty{60}{W} の電球を何年つけられますか（$c=\qty{3.00e8}{m/s}$、1年＝$\qty{3.16e7}{s}$）。
\bkpans{$\qty{9.0e12}{J}$。$9.0\times10^{12}\div60\div(3.16\times10^7)$ で、約 $4.7\times10^3$ 年（約4700年）}
\bkproblemgroup{チャレンジ}{5}
\bkproblem 地面から見て $0.8c$ で飛ぶ宇宙船について、（1）$\gamma$、（2）地面の時計で1年の間に船の中の時計が進む時間、（3）船の中ではかると \qty{100}{m} の船の、地面から見た長さを求めましょう。
\bkpans{（1）$1/\sqrt{1-0.64}=5/3\approx1.67$　（2）$1\div(5/3)=0.6$ 年　（3）$100\div(5/3)=\qty{60}{m}$}
\bkproblem 水素の \qty{486.1}{nm} の光が、ある星からは \qty{481.2}{nm} に見えました。星は近づいていますか、遠ざかっていますか。速さを求めましょう（近づくとき $\lambda=\lambda_0\sqrt{(1-\beta)/(1+\beta)}$）。
\bkpans{波長が短いので近づいている。$(481.2/486.1)^2=(1-\beta)/(1+\beta)$ より $\beta=0.0101$、約 $\qty{3.0e6}{m/s}$}
\bkproblem 運動エネルギーが静止エネルギー（\qty{0.511}{MeV}）の3倍の陽電子が、止まった電子と出会って消えます。（1）陽電子の $\gamma$・$\beta$・$pc$ を求めましょう。（2）2本の γ線が陽電子の進む向きと反対向きに出たとき、2本のエネルギーを求め、止まった陽電子のとき（⑪の PET）と比べましょう。
\bkpans{（1）$\gamma=4$、$\beta=\sqrt{1-1/16}\approx0.968$、$pc=\sqrt{2.044^2-0.511^2}=\sqrt{15}\times0.511\approx\qty{1.98}{MeV}$　（2）全エネルギー $2.044+0.511=\qty{2.555}{MeV}$。前へ出る分を $E_1$、後ろへ出る分を $E_2$ とすると $E_1+E_2=2.555$、$E_1-E_2=1.979$ なので、$E_1\approx\qty{2.27}{MeV}$、$E_2\approx\qty{0.29}{MeV}$。止まっているときは、同じ \qty{0.511}{MeV} の2本が反対向きに出る}
\bkproblem 半経験的質量公式（⑩の係数）で、鉛208（$Z=82$）の結合エネルギーを計算し、NIST の質量から求めた値 \qty{1636.4}{MeV} と比べましょう。クーロンの項はいくらですか。
\bkpans{体積 3286.4、表面 $-642.4$、クーロン $-800.4$、対称 $-215.9$、対 $+0.8$ で、約 \qty{1628}{MeV}。本当の値より約 \qty{8}{MeV} 小さい（鉛208 は陽子82・中性子126 の魔法数）。クーロンの項は約 $-\qty{800}{MeV}$}

\begin{furikaeri}
\begin{checklist}
\item 原子番号・質量数・同位体ということばで、原子核のなかみを説明できる。\bkfurref{ph:原-11:sec3}
\item 原子の質量の表から、結合エネルギーと核反応の Q 値を計算できる。\bkfurref{ph:原-11:sec4}
\item 光の時計から時間のおくれを導き、GPS の時計のずれを計算できる。\bkfurref{ph:原-11:sec7}
\item $E^2=(pc)^2+(mc^2)^2$ から、$E=mc^2$ とコンプトンの式を導ける。\bkfurref{ph:原-11:sec8}
\item 半経験的質量公式で、結合エネルギーの曲線の形を説明できる。\bkfurref{ph:原-11:sec10}
\end{checklist}
\end{furikaeri}
\bkfukidashi{mina}{太陽が光っているのは、水素が燃えているんじゃなくて、原子核がくっついて、ほんの少し軽くなったぶんがエネルギーになっていたんだね。$E=mc^2$ のおかげで、太陽は何十億年も光っていられるんだ！}
